% =====================================================================
%  Capitolo 2 — Che cos'è un gioco
%  Fonte: trascrizione della lezione GD2 (Game Design, primo anno)
%  Pacchetti richiesti nel preambolo del libro:
%    \usepackage[italian]{babel}, \usepackage[T1]{fontenc},
%    \usepackage[utf8]{inputenc} (solo pdfLaTeX), \usepackage{booktabs}
% =====================================================================

\chapter[Che cos'è un gioco]{Che cos'è un gioco: definizioni, componenti e lenti}
\label{cap:che-cos-e-un-gioco}

Il primo capitolo ha descritto il contesto del corso e il mestiere del game
designer. Questo secondo capitolo affronta la domanda da cui, in un modo o
nell'altro, discende tutto il resto: \textbf{che cos'è un gioco?} È una
domanda che sembra ovvia. Tutti sappiamo, a livello puramente euristico o
olistico, che cos'è un gioco: ognuno ha una lista di Steam o di console piena
di titoli, gioca a carte, a giochi di ruolo, a giochi da tavolo. Ma parte del
percorso formativo consiste proprio nel \textbf{passare dalla visione del
giocatore alla visione del creatore}. Un gioco visto dal giocatore e un gioco
visto dal designer sono cose molto differenti, anzi, a voler essere precisi,
completamente differenti; e questo cambio di \emph{forma mentis}, questo modo
diverso di guardare il mondo, si impara soltanto partendo dalle fondamenta.

Il percorso del capitolo ricalca quello della lezione, tenuta a tre voci dal
direttore del corso, Diego Ricchiuti, e dai due docenti di game design delle
sedi di Milano e Roma, Mattia Podini e Giacomo ``Jack'' Danna:
\begin{enumerate}
  \item ogni studente scrive la propria definizione di gioco, e le definizioni
        vengono analizzate una per una (sezioni~\ref{sec:esercizio-definizione}
        e~\ref{sec:definizioni-classe});
  \item si confrontano con alcune definizioni ``autorevoli'' e si arriva a una
        conclusione spiazzante: il gioco \emph{non è definibile}
        (sezioni~\ref{sec:definizioni-autorevoli} e~\ref{sec:non-definibile});
  \item se non si può definire il gioco, se ne possono però elencare le
        \textbf{componenti}: gli ``ingredienti della zuppa''
        (sezione~\ref{sec:zuppa});
  \item si studiano tre \emph{framework} che elencano questi ingredienti in
        modi diversi, quelli di Jesse Schell, di Brenda Romero e Ian
        Schreiber, e di Stéphane Bura (sezioni~\ref{sec:schell},
        \ref{sec:romero-schreiber} e~\ref{sec:bura});
  \item si introduce il concetto di \textbf{lente} e il fatto che non esiste un
        metodo unico di design (sezione~\ref{sec:lenti});
  \item si stabilisce il \textbf{vocabolario} del designer, a partire dalla
        prima parola che viene bandita dal corso (sezione~\ref{sec:vocabolario});
  \item si chiude con alcuni principi della \emph{forma mentis} del designer
        secondo Chris Bateman e con il task della settimana
        (sezioni~\ref{sec:bateman} e~\ref{sec:task-gd2}).
\end{enumerate}

% ---------------------------------------------------------------------
\section{Una premessa: il direttore del corso}
\label{sec:premessa-direttore}

\subsection{Chi è}

La lezione è aperta dal direttore del corso, che si presenta formalmente agli
studenti del primo anno (alcuni lo hanno già incontrato al colloquio di
ammissione). Dirige il corso di game design da più di dieci anni ed è uno
sviluppatore, precisamente un \textbf{game designer} e un \textbf{direttore
creativo}; fa videogiochi da diciassette anni. Tra i titoli a cui ha lavorato,
in vari ruoli e in varie parti dello sviluppo, cita:
\begin{itemize}
  \item un titolo multiplayer legato alla serie \emph{Poppy Playtime}, di cui
        ha curato la direzione, in uscita nel periodo di Halloween;
  \item \emph{A Quiet Place: The Road Ahead}, di cui è stato direttore
        creativo;
  \item \emph{Star Overdrive}, uscito di recente, e un altro titolo annunciato
        alla Gamescom;
  \item \emph{Dying Light 2} e \emph{Warhammer 40,000: Space Marine 2}.
\end{itemize}
È inoltre autore di due libri: \emph{Game Design Tools}, dedicato alle
tecniche di game design, e \emph{Nobody Wants Your Game: The Business Guide to
the Game Industry}, dedicato invece al lato commerciale del mestiere (vendere,
creare e gestire la propria azienda, capire i costi): una parte ``meno
romantica'' del lavoro, ma sempre molto utile.

Gli studenti lo incontreranno soprattutto al secondo e al terzo anno, ma da
quest'anno interviene anche al primo con alcune lezioni ``compassate'', come
questa. Il motivo dichiarato è la passione per l'insegnamento, la volontà di
aiutare gli altri a crescere. Nel presentarsi, sottolinea la massima fiducia e
il massimo rispetto per i due docenti di game design del primo anno, rispetto
che, avverte, ``scoprirete non essere così facile da ottenere''.

Il suo invito è a dare il 100\%: tra un anno, facendo un'introspezione, ci si
renderà conto della differenza nelle conoscenze, nelle capacità e anche nella
crescita personale. E a fare amicizia: se si è lì, si amano i videogiochi, e
attorno c'è molta gente con la stessa passione.

\subsection{Due consigli dal mestiere}

\paragraph{Giocate ai giochi di ruolo.} Il primo consiglio, che i docenti
danno di solito al secondo anno, è di \textbf{giocare a D\&D} o, più in
generale, ai giochi di ruolo da tavolo: formare gruppi (dal vivo o su Discord,
anche in digitale) e giocare. Non deve essere per forza \emph{Dungeons \&
Dragons}: vanno benissimo \emph{Pathfinder}, i giochi di \emph{Warhammer}, i
giochi narrativi moderni, \emph{Vampire: The Masquerade}, \emph{Tenebra},
\emph{Dragonbane} (un ottimo manuale uscito di recente) e la famiglia dei
cosiddetti \emph{Borg-like}, giochi costruiti sullo stesso sistema di
\emph{Mörk Borg}: \emph{CY\_BORG}, \emph{Black Powder \& Brimstone},
\emph{Pirate Borg} e altri, tra cui uno ambientato nelle rivoluzioni
proletarie. Il gioco di ruolo è \textbf{un modo spontaneo e naturale di
sperimentare il game design}: c'è tantissimo da imparare, e nel corso
dell'anno ci saranno anche lezioni dedicate, orientate non tanto al giocare
quanto alla creazione e alla comprensione di questi giochi.

È anche una cosa dell'industria: al DICE Summit, uno degli eventi più
importanti del settore rivolto agli sviluppatori (non al pubblico), chi viene
selezionato come \emph{master} di D\&D riceve un biglietto gratuito; secondo
quanto riferito in lezione, il biglietto costa diverse migliaia di dollari.

\paragraph{Prendetevi cura del vostro corpo.} Il secondo consiglio è di fare
\textbf{esercizio fisico}. Il direttore racconta di soffrire di una
protrusione alla colonna vertebrale, dovuta a uno sforzo eccessivo, che lo
costringe a fare lezione in piedi, ad andare in palestra almeno tre volte alla
settimana e a fare ginnastica posturale, perché altrimenti i dolori non lo
fanno dormire. Chi fa videogiochi passa moltissime ore seduto davanti a uno
schermo: per hobby, per passione, per lavoro e per l'accademia. Una giornata
tipo può consistere nel fare lezione seduti, giocare seduti, andare la sera a
giocare a carte seduti. Anche una sola volta alla settimana di attività fisica
è, nelle sue parole, ``il miglior investimento che potete fare sulla vostra
carriera'', perché la sedentarietà è un \emph{silent killer}.

\paragraph{Una parentesi su Commander.} Alla domanda su che mazzo giocherà
quella sera a \emph{Commander} (un formato di \emph{Magic: The Gathering} su
cui si tornerà nella sezione~\ref{sec:schell}), il direttore descrive un mazzo
costruito attorno a Mr.\ Negative, una carta che scambia i punti vita: chi
attacca rischia di ricevere i punti vita dello scambio, e questo crea una sorta
di \emph{dilemma del prigioniero} in cui nessuno vuole essere l'ultimo a
colpire. Il mazzo vince poco, ma genera una pressione psicologica molto
interessante. È già un primo esempio di come un designer guarda i giochi:
non solo ``vinco o perdo'', ma \emph{che cosa provano le persone al tavolo}.

% ---------------------------------------------------------------------
\section{L'esercizio: scrivere la propria definizione di gioco}
\label{sec:esercizio-definizione}

La lezione entra nel vivo con un esercizio molto semplice nella forma:
\textbf{scrivere la propria definizione di gioco}. Le regole:
\begin{itemize}
  \item si scrive su un blocco note digitale o su carta, indifferentemente;
  \item il tempo è di \textbf{cinque minuti} (Jack di solito ne dà due);
  \item \textbf{non} si consultano dizionari, enciclopedie o il \emph{Cambridge
        Dictionary}: il punto è dare la \emph{propria} definizione;
  \item una volta scaduto il tempo, la definizione \textbf{non va
        modificata} mentre si ascoltano quelle degli altri. Non c'è nessun
        premio: è una questione di condivisione e di comprensione.
\end{itemize}
Anche i tre docenti scrivono (o hanno già pronta) la loro definizione, che
verrà letta per ultima.

Le definizioni vengono poi lette una per una e analizzate. Il metodo di
analisi è sempre lo stesso, ed è importante capirlo prima di leggere la
sezione seguente: per ogni definizione si cercano \textbf{gli elementi che
include e quelli che esclude}. Ogni parola scelta lascia fuori qualcosa che
pure chiameremmo ``gioco''. I docenti lo ripetono più volte: non si tratta di
essere ``infami'' o di dimostrare che una definizione è sbagliata, ma di
mostrare come ognuna contenga \emph{un'idea} di gioco senza riuscire ad
abbracciare l'intera definizione di gioco, per quanto ampia possa essere.

Un'avvertenza che emerge subito: \textbf{le parole hanno un significato molto
preciso}. Se si scrive ``spazio'' intendendo anche uno spazio immaginario, ma
non lo si specifica, si è scritto ``spazio'', non ``spazio immaginario''. Allo
stesso modo, odiando tutti un po' la matematica, ``numeri'' e ``numeri
immaginari'' sono cose estremamente differenti: definire i numeri non
definisce i numeri immaginari.

% ---------------------------------------------------------------------
\section{Le definizioni della classe}
\label{sec:definizioni-classe}

Si riportano di seguito le definizioni lette in aula (in presenza e online),
nell'ordine in cui sono state discusse, ciascuna con il commento dei docenti.
Gli studenti sono indicati con il solo nome di battesimo.

\subsection{Un prodotto con giocatori, regole, spazio, tempo e scopo}

\begin{quote}
\emph{``Un prodotto a scopo ludico o di intrattenimento che presenta diverse
caratteristiche comuni: dei giocatori, delle regole, uno spazio e un tempo
definiti e uno scopo e/o degli obiettivi.''} (Filippo)
\end{quote}

\begin{description}
  \item[Prodotto] è un termine \textbf{commerciale}: un prodotto definisce una
        vendita. Scegliendolo si escludono giochi che chiamiamo
        tranquillamente tali: \emph{mosca cieca}, ad esempio, non è un
        prodotto e non può essere venduto. È sbagliato? No: è una definizione.
  \item[Scopo ludico] implica che un gioco abbia una funzione di
        intrattenimento. Ma esistono giochi il cui scopo non è affatto ludico.
        Il caso esemplare è \emph{That Dragon, Cancer}: un gioco creato da uno
        sviluppatore che racconta la propria vera storia, quella di un padre
        che deve accettare che il figlio di quattro anni morirà di cancro. È
        straziante, ``una delle cose più brutte che spero non mi accada mai'',
        e giocarlo non ha nulla di ludico; eppure dura circa due ore, è
        altamente consigliato, è su Steam ed è anche un prodotto. È un gioco.
  \item[Spazio] (osservazione di Jack) esistono giochi che si svolgono in uno
        spazio puramente immaginario: dov'è lo spazio di gioco degli
        indovinelli tra Bilbo e Gollum nello \emph{Hobbit}? Lo studente
        intendeva lo spazio anche in senso non fisico, ma non lo ha scritto:
        ``spazio'' e ``spazio immaginario'' sono cose molto differenti.
\end{description}

\subsection{Un'attività tra persone con regole concordate}

\begin{quote}
\emph{``Attività tra persone che, tramite regole concordate, passano il tempo
insieme con il fine di essere intrattenute.''} (Damiano)
\end{quote}

\begin{description}
  \item[Attività] indica una qualche forma di azione. Esistono però giochi
        quasi privi di attività: gli \emph{autoplay} e gli \emph{idle} (o
        AFK), in cui si prende qualche decisione all'inizio e poi il gioco va
        avanti da sé, o addirittura giochi che ``si autogiocano''. Su Twitch
        esistono streamer che fanno giocare animali: un pesce rosso ha finito
        \emph{Pokémon Rosso} (impiegando l'equivalente di centinaia di giorni)
        e un pesce rosso ha sconfitto Malenia in \emph{Elden Ring}. Il
        direttore, che ci ha messo ventiquattro tentativi, si chiede quanti ne
        abbia impiegati il pesce.
  \item[Regole concordate] i bambini piccoli giocano molto spesso senza
        condividere le regole: ``siamo pirati, però adesso arriva un mostro,
        però adesso divento un Super Saiyan, però dovevamo partire da
        dietro''. Il gioco è un cambiamento costante, ``melassoso'', senza un
        inizio e senza una fine, e le regole non sono mai state concordate.
\end{description}

\subsection{Un'attività o un oggetto che stimola emozioni e pensiero}

\begin{quote}
\emph{``Un'attività o un oggetto capace di intrattenere il giocatore
stimolando emozioni e pensiero.''} (Federico)
\end{quote}

\begin{description}
  \item[Oggetto] è la prima volta che compaiono gli oggetti. Non è sbagliato,
        ma quando si parla di oggetti spesso si parla di \textbf{giocattoli},
        e come si vedrà più avanti nel corso c'è una grossa differenza tra un
        giocattolo e un gioco; ciò non toglie che esistano giochi che sono
        oggetti.
  \item[Stimolare emozioni e pensiero] esistono giochi molto passivi, che
        funzionano in modo estremamente automatico, ``che si fanno con metà
        del cervello spenta, se non tutto''. \emph{Watch Paint Dry}, che ha
        venduto migliaia di copie su Steam, è una schermata monocolore che
        diventa leggermente più chiara col passare del tempo, perché simula il
        guardare una parete che si asciuga. \emph{Placid Plastic Duck
        Simulator}, un successo italiano realizzato da un docente
        dell'accademia, ha venduto centinaia di migliaia di copie e fatto
        milioni di euro: è un gioco in cui si guardano delle paperelle di gomma
        muoversi in un ambiente.
\end{description}

\subsection{Un prodotto che regala stimoli ed emozioni}

\begin{quote}
\emph{``Un prodotto che regala stimoli ed emozioni a chi gioca e in parte anche
a chi sviluppa.''} (Matteo)
\end{quote}

Torna il termine \emph{prodotto}, con il suo fine commerciale, che esclude i
giochi di parole, i giochi infantili e uno spettro molto ampio di altre
attività. Torna soprattutto il discorso delle \textbf{emozioni}, che inizia a
essere una costante: il cosiddetto \textbf{responso emotivo} è un punto molto
importante (è presente anche nella definizione di Mattia) e nel corso degli
anni si vedranno diversi modi per definirlo. Restano però fuori elementi come
lo spazio, il tempo, la condivisione e il fatto che il gioco sia o non sia
giocato da un essere umano. Se a giocare a scacchi è una macchina (come i
sistemi di DeepMind, \emph{AlphaGo} per il go e \emph{AlphaZero} per gli
scacchi), e magari contro un'altra macchina, nessuno prova emozioni: eppure il
gioco esiste. Esiste ``nel vuoto''; ma una volta messo sul mercato il discorso
cambia, e si pone una domanda che tornerà: \textbf{se un gioco esiste ma
nessuno lo gioca, che valore ha?} È la versione ludica della classica domanda
filosofica: se un albero cade nella foresta e nessuno lo sente, fa rumore?

\subsection{Un'esperienza di immedesimazione e di fuga dalla realtà}

\begin{quote}
\emph{``Un gioco è un'esperienza in grado di far immedesimare al meglio il
giocatore, suscitare emozioni e sentimenti, far sviluppare senso critico e
immaginazione: un modo per fuggire dalla realtà.''} (Sara)
\end{quote}

È all'opposto della precedente: invece che molto ampia, è estremamente
stretta.
\begin{itemize}
  \item \textbf{Immedesimazione}: i giochi che non prevedono di immedesimarsi
        non sarebbero giochi.
  \item \textbf{Senso critico}: il \emph{gioco dell'oca} non fa sviluppare
        alcun senso critico (nella migliore delle ipotesi è un passatempo);
        \emph{Cookie Clicker} nemmeno. Secondo questa definizione non sarebbero
        giochi, come migliaia di altri titoli.
  \item \textbf{Fuggire dalla realtà}: è il punto più interessante. Una delle
        caratteristiche dei videogiochi (non di tutti, e non dei giochi in
        generale) è l'immaginazione, il poter creare cose che non esistono nel
        mondo reale: c'è \emph{Modern Warfare 2} con le armi, ma c'è anche
        \emph{Skyrim} con i draghi, la fantascienza, \emph{Gothic} con gli
        orchi (``per fortuna gli orchi non esistono''). Non a caso si parla di
        \emph{escape from reality}: non solo inventare o immaginare, ma
        \emph{fuggire} dal mondo reale, che è complesso, difficile e spesso
        brutto (basta guardare le guerre in corso o millenni di storia umana).
        Il videogioco può presentare qualcosa di molto diverso, unico, magico e
        potenzialmente bellissimo.
\end{itemize}

\subsection{Emozioni positive e negative}

Una definizione successiva (Gaetano) torna sull'\textbf{attività}, sulle
\textbf{emozioni}, sull'\textbf{interattività}, ma aggiunge un elemento nuovo:
le emozioni possono essere \textbf{positive e negative}. In questo modo
\emph{That Dragon, Cancer} rientra nella definizione; restano invece fuori i
prodotti privi di interattività. E qui nasce una domanda legittima: se un
gioco può non essere interattivo, allora anche un film può essere un gioco? La
risposta è: \emph{potenzialmente sì}, e ci si arriverà (sezione~\ref{sec:zuppa}).
Restano esclusi anche i giochi che non danno emozioni, come \emph{Watch Paint
Dry}.

Il concetto di emozione negativa è molto interessante: certi giochi si
giocano \emph{per soffrire}, per spaventarsi o per arrabbiarsi. Contro Malenia
il direttore ha perso ventiquattro volte: era frustrato e arrabbiato, il
controller ``ha sofferto'', eppure ha continuato a giocare. Quando l'ha
sconfitta, il grido gutturale di gioia avrà fatto pensare ai vicini di
chiamare il 112; i tentativi precedenti avranno prodotto grida di odio, rabbia
e sofferenza.

\subsection{Un'esperienza individuale o di gruppo, con un obiettivo}

\begin{quote}
\emph{``Un'attività o un'esperienza individuale o di gruppo con regole e
meccaniche prestabilite e un obiettivo da raggiungere, a volte fine
all'esperienza stessa, con il fine di intrattenere e far scaturire differenti
emozioni al giocatore.''} (Francesca)
\end{quote}

\begin{description}
  \item[Individuale o di gruppo] elemento non ancora emerso: esistono il
        \emph{single player} e il \emph{multiplayer}, giochi digitali e non
        digitali, di gruppo e non di gruppo. Il videogioco può essere
        un'esperienza solitaria o un'esperienza condivisa.
  \item[Obiettivo] esistono giochi senza obiettivo. In \emph{Proteus} si
        cammina in un campo fiorito, e basta: nessun obiettivo, nessun inizio,
        nessuna fine; a un certo punto, si suppone, ci si annoia e lo si
        chiude. Un gioco ``molto carino'', ma senza obiettivo.
  \item[A volte fine all'esperienza stessa] è la parte più interessante:
        l'obiettivo può essere \emph{interno} al gioco (sconfiggi il cattivo,
        salva o conquista il mondo, distruggi l'impero) oppure \emph{non
        legato al gioco stesso}. In \emph{Magic: The Gathering} lo scopo delle
        carte sarebbe giocare a Magic, ma esiste un intero gruppo di persone il
        cui obiettivo è \textbf{collezionare} le carte, cosa che non ha nulla
        a che vedere con le partite, eppure fa parte dell'esperienza di Magic.
        In altri giochi questa componente è addirittura quella principale: il
        gioco di carte \emph{Pokémon} soffre in questo periodo di un'eccessiva
        collezionabilità.
\end{description}
Anche qui, nonostante la definizione sia relativamente ampia, l'aver inserito
gli obiettivi esclude i giochi privi di obiettivi.

\subsection{Un'attività intrattenente basata sull'interattività}

\begin{quote}
\emph{``Un'attività intrattenente che si basa sull'interattività di colui che
ne usufruisce.''} (Flavio)
\end{quote}

\begin{description}
  \item[Intrattenente] esce dal campo delle emozioni: intrattenere significa
        anche, semplicemente, \textbf{non annoiare}. L'assenza di noia è già
        una forma di intrattenimento. È il motivo del successo, qualche anno
        fa, dei \emph{fidget spinner}: non sono piacevoli in sé, ma rispetto
        alla noia sono meglio. Molte cose sono meglio della noia, che è una
        delle sensazioni peggiori che il cervello possa vivere.
  \item[L'interattività di chi ne usufruisce] esclude chi non sta giocando.
        Ma esistono giochi interattivi \emph{anche per chi non gioca}: i giochi
        pensati per gli streamer. \emph{A Quiet Place} ha una modalità
        streaming in cui gli spettatori possono ``incasinare'' lo streamer:
        stanno interagendo, ma non sono i fruitori diretti. Sono un
        \textbf{secondo livello di fruitori}.
\end{description}

\subsection{Vivere una storia in prima persona}

\begin{quote}
\emph{``Un gioco è un'esperienza interattiva in cui il giocatore vive in prima
persona una storia, interpreta un ruolo, compete o collabora con altri
giocatori per raggiungere un obiettivo, immergendosi nel contesto creato.''}
(Francesco)
\end{quote}

\begin{description}
  \item[In prima persona] non va inteso nel senso della camera, ma nel senso
        di impersonare un \emph{avatar} (o $n$ avatar). Esistono anche giochi
        che rendono il rapporto molto più contorto, mettendo il giocatore nei
        panni di una camera che osserva l'avatar.
  \item[Immergersi nel contesto creato] esistono giochi
        \textbf{metanarrativi} che separano deliberatamente il giocatore
        dall'avatar, cercando una \emph{dissonanza ludonarrativa} tra chi gioca
        e ciò che viene giocato: \emph{Inscryption}, \emph{The Stanley
        Parable}, \emph{The Beginner's Guide} (dello stesso autore di
        \emph{Stanley Parable}).
  \item[Storia] la parola, per quanto ovvia, non era ancora stata usata. Ci
        sono giochi che sono praticamente storie, come \emph{Gone Home} o in
        generale i cosiddetti \emph{walking simulator}, e giochi che sono puro
        gameplay, come il flipper o \emph{Tetris}. Una definizione centrata
        sulla storia esclude tutti i giochi puramente meccanici.
  \item[Contesto] esistono giochi in cui il contesto non esiste e viene creato
        dai giocatori stessi. \emph{Ten Candles} è un gioco di ruolo che si
        gioca con dieci candele (da cui il nome), e parte del gioco è proprio
        creare il gioco stesso: il contesto nasce da zero. Capita in moltissimi
        \textbf{giochi di ruolo di stampo narrativo} (in Italia ne pubblica
        molti, ad esempio, Narrattiva). Tra i preferiti del direttore c'è
        \emph{Passione delle Passioni}, che simula una soap opera alla
        \emph{Beautiful}: i personaggi non muoiono mai davvero, e quando il
        proprio personaggio muore torna due turni dopo con un colpo di scena
        (``era mio cugino'', ``in realtà sono naufragato su un'isola deserta e
        sono sopravvissuto per due anni''). Jack aggiunge che la genialità del
        gioco sta nel chiedere di interpretare \emph{sia l'attore sia il
        personaggio} che l'attore interpreta.
\end{description}

\subsection{Un'avventura da vivere come un alter ego}

\begin{quote}
\emph{``Un'avventura in cui il giocatore vuole immergersi e viverla il più
possibile, come se fosse una sorta di alter ego in cui vivi quel mondo nelle
vesti del protagonista.''} (Lorenzo)
\end{quote}

Una definizione molto propositiva. La parola \textbf{avventura} propone già
l'idea di qualcosa di più interessante del mondo reale, che è decisamente meno
avventuroso. C'è poi un elemento di immersione e soprattutto di
\textbf{volontà di partecipazione}, implicito ma percepibile, che nessuno aveva
ancora citato: la \textbf{volontarietà}. Un gioco è un'esperienza volontaria?
Dipende: ci sono giochi volontari e giochi non volontari. Con la sola parola
``avventura'', ovviamente, si potrebbe escludere di tutto.

\subsection{L'interazione con un medium, con l'ambiente o con la propria immaginazione}

\begin{quote}
\emph{``L'interazione attiva con un medium, con l'ambiente circostante o con la
propria immaginazione, il cui scopo è principalmente quello di divertirsi e
farsi coinvolgere.''} (Federica)
\end{quote}

La studentessa spiega di aver voluto essere molto generica, per includere
molte cose: distinguere il gioco dal guardare un film, ma includere anche i
giochi della sua infanzia, come ``giocare a pensare'', immaginare la vita di
una persona, o saltare la campana. ``Giocare a pensare'' è senza dubbio un
gioco, e qui si apre un intero mondo non ancora trattato: il
\textbf{gioco pensato}. Non solo lo spazio immaginario, ma l'interazione
immaginaria e le sue conseguenze: l'intera esperienza esiste in un ambito
completamente fluido. È ciò che accade ai bambini che giocano insieme pur con
regole completamente diverse, mai concordate.

Esiste un'intera branca della matematica dedicata a questi giochi: i
\textbf{puzzle di logica}, interamente mentali, in cui si devono individuare
o capire le regole. Oppure gli indovinelli del tipo ``sei su un tetto, come
fuggi?'', in cui si possono fare solo domande a cui si risponde sì o no, o in
cui si deve ricostruire come si è arrivati in una certa situazione: giochi
deduttivi, ambientati in uno spazio puramente mentale, che si possono fare
anche da soli. In ambito scientifico si parla di \textbf{esperimenti
mentali}. Uno dei più famosi è il \emph{paradosso dei gemelli} legato alla
relatività di Einstein: se uno di due gemelli identici viaggia nello spazio a
velocità prossime a quella della luce e l'altro resta sulla Terra, al ritorno
saranno ancora identici? La risposta è no. Einstein non ha mai spedito due
gemelli nello spazio: era un esperimento mentale.\footnote{L'effetto di
dilatazione temporale previsto dalla relatività è stato poi verificato
sperimentalmente, ad esempio con orologi atomici trasportati in volo, ma
l'esperimento dei gemelli in quanto tale resta un esperimento mentale.}

Jack fa notare che è stata anche la prima persona a usare, al posto di
``intrattenimento'', il termine \emph{divertimento}: una parola su cui si
tornerà prima della fine della lezione (sezione~\ref{sec:vocabolario}).

\subsection{Un insieme di regole e sistemi}

\begin{quote}
\emph{``Il gioco è un insieme di regole e sistemi con uno scopo ben
preciso.''} (Patrick)
\end{quote}

Compare per la prima volta la parola \textbf{sistemi}. Finora si era parlato
di regole e obiettivi; i sistemi aggiungono un livello ulteriore, perché le
regole si incastrano in strutture più complesse delle regole stesse.

\paragraph{Esempio: il contropiede.} Nel regolamento del calcio non esiste la
``regola del contropiede''. Il contropiede è una tattica che \textbf{deriva da
come le regole si incastrano}: qualcuno ha capito che, se si lancia la palla
dall'altra parte del campo quando gli avversari sono tutti nella propria metà
e si corre ``come pazzi'', si ottiene un vantaggio strategico. Nessuna regola
lo dice: è un sistema che nasce naturalmente dal gioco. Allo stesso modo,
chiunque abbia scoperto una \emph{build} ``rotta'' in un gioco ha trovato
qualcosa che nessuna regola descriveva (``se fai A più B rompi il gioco''):
lo ha scoperto analizzando i sistemi.

Si apre così un intero mondo di cui parlare: che cos'è un sistema? Quando
nasce? Che cosa rende un gioco sistemico o non sistemico (si veda anche il
capitolo~\ref{cap:introduzione-game-design})? Dopo una quindicina di
definizioni, con altre nove ancora da leggere, nessuno aveva ancora toccato
questo punto.

\subsection{Un insieme di regole da rispettare}

\begin{quote}
\emph{``Un gioco è un'attività definita da un insieme di regole da rispettare al
fine di raggiungere un obiettivo.''} (Davide)
\end{quote}

Tornano obiettivo e regole, ma l'espressione ``da rispettare'' porta un
elemento nuovo. Quanti hanno \textbf{barato} almeno una volta? A Monopoly,
prendendo una banconota da cento quando nessuno guarda; a \emph{UNO}, mettendo
un $+2$ su un altro $+2$, che tecnicamente non è nelle regole ufficiali. Ci
sono giochi che addirittura incorporano l'imbroglio: \emph{Munchkin} ha una
carta che si chiama proprio ``Barare!'' e dà il potere di farlo. E anche negli
sport si bara: in Formula 1, quasi ogni anno, c'è uno scandalo per qualcuno
che ha manomesso il motore o ``sistemato'' qualche componente. Si è barato, e
si è comunque giocato.

Se con Lorenzo era emersa la volontarietà, con Davide emerge
l'\textbf{onestà}, il tentativo di rispettare le regole: è vero che si bara,
ma se ognuno facesse costantemente ciò che vuole il gioco sarebbe puro caos.

\subsection{Un'attività interessante e non obbligatoria}

\begin{quote}
\emph{``Un gioco è un'attività interessante, non obbligatoria, per passare il
tempo.''} (Federico)
\end{quote}

\begin{itemize}
  \item \textbf{Interessante} esclude la distinzione tra positivo e negativo:
        torna l'idea che ``piuttosto che niente, è meglio piuttosto''.
  \item \textbf{Non obbligatoria} ripropone la volontarietà.
  \item \textbf{Per passare il tempo} esclude tutti i giochi che hanno una
        qualche valenza di crescita: ad esempio i primi \emph{Assassin's
        Creed}, diventati famosi perché in grado di insegnare la storia, a
        detta di alcuni, meglio dei libri di storia. Ma introduce anche un
        elemento che nessuno aveva considerato: i \textbf{giochi passatempo}.
        Sono più semplici, più banali se si vuole, ma esistono, e la loro
        funzione è proprio essere semplici. Lo stesso direttore ha sul
        telefono giochi che non richiedono grande impegno né grande
        \textbf{carico cognitivo}, adatti ai due o tre minuti di attesa
        dell'autobus o del treno.
\end{itemize}

\subsection{Un regolamento creato da un designer}

\begin{quote}
\emph{``Il gioco è un modo di svagarsi sottostando a un regolamento''},
a cui lo studente aggiunge: \emph{``\dots di un designer''}. (Francesco)
\end{quote}

\emph{Svagarsi} richiama il passare del tempo, non necessariamente in modo
positivo. La novità è l'aggiunta del \textbf{designer}: le regole sono state
progettate da qualcuno. Questo, però, esclude il calcio, che è un
\textbf{risultato evolutivo}: non esiste un inventore del calcio, che ha
antenati antichissimi (si citano i giochi con la palla delle civiltà
mesoamericane) e si è evoluto nei secoli fino alla forma moderna.

Anche ammettendo un designer, chi è? Il direttore è stato direttore creativo
di \emph{A Quiet Place}, ``la punta di diamante del design'', ma non l'unico
designer: il design nasce da una visione di idee e di sistemi che coinvolge
più persone. Si potrebbe allora dire ``un gruppo di creatori''; ma che cos'è un
gruppo, quanto è grande, come lo si definisce? All'estremo opposto ci sono i
giochi di Hideo Kojima, che mette il proprio nome ovunque nei titoli di
\emph{Death Stranding} (``A Hideo Kojima Game'', ``Directed by Hideo
Kojima''\dots): nei primi trenta minuti la parola ``Kojima'' compare un numero
impressionante di volte.

\subsection{Interazioni che influenzano il gioco, sfida e progressione narrativa}

\begin{quote}
\emph{``Un qualunque medium di intrattenimento con interazioni che hanno
un'influenza su di esso, con un livello di sfida e/o una progressione
narrativa.''} (Eric)
\end{quote}

\begin{description}
  \item[Progressione narrativa] una storia che si evolve. Una storia può
        essere un semplice punto di partenza; qui invece la storia
        \emph{progredisce}. Attenzione: ``progredisce'' non implica un
        finale.
  \item[Interazioni che hanno effetto sul gioco] esclude le interazioni che
        non hanno effetto sul sistema, ad esempio quelle \textbf{sociali}: in
        un gioco multiplayer ridere e scherzare con chi sta giocando è
        un'interazione, ma non modifica il sistema di gioco.
\end{description}

\subsection{Un'esperienza attraverso uno schermo, fatta di personaggi}

\begin{quote}
\emph{``Un gioco è un'esperienza, qualcosa che riesce a darti una determinata
soddisfazione attraverso uno schermo; è solitamente fatto di storia,
personaggi, luoghi, ambientazioni e tanto altro che ti fanno provare certe
sensazioni.''} (Francesco)
\end{quote}

\emph{Attraverso uno schermo} esclude tutto ciò che non è videogioco o gioco
mobile. Ma compaiono per la prima volta i \textbf{personaggi}: avatar del
giocatore, personaggi, NPC sono elementi estremamente comuni nei giochi e
fino a quel momento nessuno li aveva citati direttamente. Esistono giochi
senza personaggi? Moltissimi. Esistono giochi basati \emph{interamente} sui
personaggi, fatti solo di interazioni con loro? Sì: uno verrà analizzato in
una lezione specifica, e non viene nominato per non rovinare la sorpresa.

\subsection{Un'opera d'arte}

\begin{quote}
\emph{``Un gioco è un'opera d'arte, un modo evoluto di raccontare un mondo
facendolo vivere in prima persona.''} (Mattia)
\end{quote}

Lo studente precisa che ``raccontare un mondo'' può avvenire anche attraverso
\emph{Pong}, attraverso la simulazione di una persona che gioca a ping pong, o
può riguardare la storia stessa dello sviluppo del gioco. L'elemento chiave,
però, è l'\textbf{opera d'arte}, concetto spesso in contrapposizione con
quello di \emph{prodotto} emerso nelle prime definizioni.

Esistono tante opere e poche opere d'arte: di violinisti, compositori e
pianisti è pieno il mondo passato e presente, ma tra secoli ricorderemo ancora
Mozart. Che cos'è allora un'opera d'arte, e chi la definisce? Una risposta
classica è ``ai posteri l'ardua sentenza''. Un esempio caro al direttore è
\emph{Moby Dick}: alla sua uscita il romanzo fu accolto in modo pessimo, fu
dimenticato per decenni e venne riscoperto solo molto tempo dopo, quando i
critici lo rilanciarono come una delle più grandi opere mai scritte; oggi è
considerato da tutti un capolavoro.\footnote{In lezione l'episodio è
raccontato in forma aneddotica (copie bruciate, duecento anni di oblio).
Storicamente il romanzo di Herman Melville (1851) ebbe un'accoglienza tiepida
e vendite scarse, parte delle copie andò distrutta nell'incendio di un
magazzino dell'editore nel 1853, e la sua riscoperta avvenne con la cosiddetta
\emph{Melville Revival} degli anni Venti del Novecento.} Non sappiamo quindi
come si definisca un'opera d'arte; ma questa definizione aggiunge un
\textbf{elemento romantico}: il mirare a creare qualcosa di grande, che ci si
riesca o meno, una volontà di elevazione.

\subsection{Divertire, insegnare, far riflettere, meravigliare}

\begin{quote}
\emph{``Un gioco è un'attività di intrattenimento che ha lo scopo di divertire
ma può anche insegnare, far riflettere o meravigliare, e fa tutto ciò in uno
spazio di regole e intenzioni definiti.''} (Claudio)
\end{quote}

Tornano varie emozioni, tornano le regole e torna lo \textbf{spazio},
stavolta \emph{definito}, cioè in qualche modo percepibile. Con i bambini non
funzionerebbe, ma ci sono giochi in cui la definizione dello spazio è il
punto di forza: \emph{Pokémon GO} sovrappone uno spazio digitale allo spazio
del mondo reale, e così faceva il suo ``papà'', \emph{Ingress}.

\subsection{Un'attività per la persona giocante}

\begin{quote}
\emph{``Un'attività interattiva svolta per divertire o intrattenere la persona
giocante.''} (Antonino)
\end{quote}

L'elemento nuovo è la \textbf{persona giocante}. Molte delle definizioni
ascoltate \emph{escludono il giocatore}: non lo nominano come parte della
definizione. Si torna così alla domanda posta da Mattia: esiste un gioco senza
giocatori? È di nuovo la domanda dell'albero che cade nella foresta, una
questione filosofica antica.

\subsection{Un'attività per una o più persone}

\begin{quote}
\emph{``Un gioco è un'attività volta al divertimento, che può essere di diverse
tipologie e comprendere più persone.''} (Valerio)
\end{quote}

Molto generica. \emph{Volta al divertimento} esclude i giochi tristi; ritorna
invece, rispetto alla precedente, il tema del \textbf{gruppo}: uno o più
giocatori.

\subsection{Una promessa e un punto di vista diverso dal proprio}

\begin{quote}
\emph{``Un gioco promette [\dots]: permette al giocatore di prendere i panni del
protagonista della storia per fargli vivere un'esperienza unica nel suo
genere, facendogli vedere la realtà sotto il punto di vista dei personaggi e
consentendogli di far valere le sue scelte morali.''} (Andrea)
\end{quote}

Lo studente premette di aver già capito, ascoltando le altre, che nella sua
c'è qualcosa che non va. ``Ma in tutte!'', è la risposta: non ne è stata
trovata nemmeno una senza buchi.

\paragraph{La promessa.} \emph{Promettere} è un termine interessantissimo.
Quante volte ci si è arrabbiati con un gioco o con un film perché il trailer
era diverso dall'esperienza? Un gioco, in un certo senso, deve definire se
stesso. Non sempre lo fa: le regole del Monopoly non dicono come si
\emph{vive} il Monopoly, e il regolamento del calcio non descrive le emozioni
del calcio. Altri giochi invece sì: un \emph{God of War} cerca di spiegare
come si giocherà e quale sarà l'esperienza; i libretti di istruzioni dei
giochi di qualche generazione fa avevano proprio questa funzione. La promessa
vale anche in altri ambiti: un ristorante che si chiama ``Sushi King'' e poi è
una braceria senza sushi è strano, indipendentemente da quanto sia buona la
brace.

\paragraph{Il punto di vista dei personaggi e le scelte morali.} Il gioco
permette di vivere condizioni che non corrispondono necessariamente alla
propria attitudine, ad esempio \textbf{interpretare il cattivo}. In
\emph{Fallout 3} c'è una scelta morale in cui un personaggio (Mr.\ Burke)
offre al giocatore una suite di lusso nella Tenpenny Tower in cambio della
detonazione della bomba atomica al centro della città di Megaton. Il
direttore e uno studente l'hanno fatta esplodere; un altro studente,
``sofferente'', ha rinunciato. Nella vita reale nessuno farebbe esplodere una
città per una casa; nel gioco sì, perché si sta vivendo l'esperienza di un
personaggio cattivo.

In \emph{Baldur's Gate 3} una delle \emph{run} possibili è quella malvagia,
con scelte moralmente discutibili. Jack racconta di trovarla straziante: già
una delle prime scelte, saccheggiare il Boschetto di Smeraldo invece di
proteggerlo e aprire il cancello ai goblin, porta a una scena in cui Zevlor
guarda il giocatore non con rabbia, ma con l'incredulità di chi chiede ``perché
ci stai condannando tutti?''. Mantenere una run malvagia è difficile, perché
moralmente non si farebbe mai.

Il punto nuovo, quindi, non è la storia o la scelta (già emerse), ma
l'\textbf{interpretazione di punti di vista potenzialmente non propri}, nel
bene e nel male.

\subsection{Creare una specifica esperienza}

\begin{quote}
\emph{``Un gioco è un insieme di regole e sistemi con l'obiettivo di creare una
specifica esperienza e far provare delle emozioni ai giocatori.''} (Tancredi)
\end{quote}

L'inizio è simile a quello di Patrick; il termine più importante è
\textbf{specifica esperienza}. C'è una volontarietà simile a quella del
designer citata da Francesco: voler creare qualcosa di specifico, di meno
casuale. Eppure esistono moltissimi giochi casuali: la \textbf{roulette} è un
gioco puramente casuale (se non si bara).

\subsection{Scenari, meccaniche e personaggi ideati dal designer}

\begin{quote}
\emph{``Un gioco è un'esperienza attraverso la quale il giocatore, ovvero chi
lo gioca, riesce a immedesimarsi e/o a provare emozioni che il designer ha
ideato con l'introduzione o la costruzione di scenari, meccaniche e
personaggi.''} (Angelo)
\end{quote}

Compaiono un termine tecnico, \textbf{meccaniche}, che verrà studiato a
fondo, e il concetto di \textbf{scenario}. Uno scenario indica che il
contesto, sia fisico sia di spazio (digitale o meno), è stato creato: quando si
immagina uno scenario si immaginano tutti i suoi elementi. Nessuno aveva
ancora parlato di scenari, che pure sono onnipresenti nei giochi: anche un
puzzle è uno scenario, un problema e la sua soluzione.

\subsection{Un'esperienza sempre diversa}

\begin{quote}
\emph{``È un medium attraverso il quale si intrattiene il giocatore, con
scelte e obiettivi, e quindi ogni volta l'esperienza del giocatore è
differente.''} (Eric)
\end{quote}

\begin{itemize}
  \item \textbf{Ogni volta differente} esclude moltissimo: esistono giochi
        puramente lineari in cui l'esperienza è sempre identica.
  \item \textbf{Si intrattiene il giocatore} esclude chi non gioca: ci sono
        giochi che si giocano e giochi che si \emph{guardano giocare}. Il
        direttore guarda giochi in \emph{blind run} senza commento, proprio
        per godersi il gioco e non lo streamer. Un gioco può intrattenere
        anche chi non lo sta giocando.
\end{itemize}

\subsection{Un viaggio o una sfida}

\begin{quote}
\emph{``Un gioco può essere un viaggio piacevole o una sfida, che può essere
in compagnia o in solitaria.''} (Alberto Maria)
\end{quote}

La parola nuova è \textbf{sfida}. Nessuno aveva ancora detto esplicitamente che
un gioco può essere una \emph{challenge} da superare, qualcosa di complesso:
Malenia era stata citata più volte, ma il concetto di sfida era rimasto
implicito.

% ---------------------------------------------------------------------
\section{Le definizioni ``autorevoli''}
\label{sec:definizioni-autorevoli}

Dopo circa un'ora di definizioni ``tutte sbagliate'', ci si potrebbe aspettare
la definizione giusta. Ne vengono proposte alcune, con un avvertimento:
\emph{saranno comunque sbagliate}.

\begin{description}
  \item[Oxford Dictionary] un gioco è un'attività in cui ci si impegna
        (\emph{engage}) per divertimento (\emph{amusement} o \emph{fun}), come
        forma di competizione o di sport, giocata seguendo delle regole.
  \item[Friedrich Schiller] il \emph{play} è ``l'espressione priva di scopo
        di un'energia esuberante'' (\emph{the aimless expenditure of
        exuberant energy}). Qui serve una precisazione linguistica: in
        italiano ``gioco'' indica sia lo strumento sia l'attività, mentre
        l'inglese distingue \textbf{game} (il gioco come oggetto, struttura) e
        \textbf{play} (l'atto di giocare). È uno dei rari casi in cui l'inglese
        è più preciso dell'italiano. La definizione di Schiller è l'opposto di
        molte di quelle della classe: è un'attività \emph{priva di obiettivo},
        e non serve a passare il tempo, ma a spendere un'energia in eccesso,
        quando in un certo momento se ne ha più di quanta ne serva, senza
        alcuna ragione.
  \item[Tracy Fullerton, Chris Swain e Steven Hoffman] (autori di \emph{Game
        Design Workshop}) un gioco è un \textbf{sistema formale} (cioè
        definito, chiuso) in cui i giocatori interagiscono in una
        \textbf{struttura di conflitto} che si risolve in un \textbf{esito
        ineguale}. \emph{Conflitto} non significa guerra, ma semplicemente
        un'interazione che porta a un risultato in cui qualcuno vince e
        qualcuno perde, qualcuno tiene i soldi e qualcuno no. Ovviamente
        esclude i giochi in cui vincono tutti, come i \textbf{cooperativi}.
  \item[Una citazione sul play] riportata in lezione: ``la verità riguardo al
        \emph{play} è che non può essere definito, così come non può essere
        definita la vita stessa''.
\end{description}

% ---------------------------------------------------------------------
\section{Il gioco non è definibile}
\label{sec:non-definibile}

\subsection{Il gioco come stato mentale}

L'ultima citazione è l'intero punto dell'esercizio:
\begin{quote}
\centering
\textbf{Nessuna delle vostre definizioni è giusta. Nessuna di queste
definizioni è giusta. Nessuna delle vostre definizioni è sbagliata.}
\end{quote}
Il gioco \textbf{non è definibile}, perché alla base il gioco è
\textbf{uno stato mentale}: qualcosa che si interpreta. Si potrebbe decidere
di ``giocare'' il proprio lavoro, di interpretare la propria vita come un
gioco. Lo dice bene una citazione da \emph{Mary Poppins}:
\begin{quote}
\emph{In every job that must be done there is an element of fun. You find the
fun and --- snap! --- the job's a game.}
\end{quote}
(In ogni lavoro che va fatto c'è un elemento di divertimento: trovalo e,
schiocco di dita, il lavoro diventa un gioco.) C'è anche un detto italiano
simile: ``scegli un lavoro che ami e non lavorerai un solo giorno della tua
vita''. Spoiler: non è così, perché il lavoro resta lavoro; ma è meglio
lavorare su qualcosa che si ama che su qualcosa che si odia.

Tuttavia, se il gioco non è definibile, è definibile \emph{ciò che
caratterizza} un gioco (sezione~\ref{sec:zuppa}).

\subsection{Le definizioni dei docenti}

Prima di procedere, i tre docenti leggono le proprie definizioni, con lo
stesso avvertimento: non vanno prese come ``quella giusta''. Ognuna esclude una
marea di giochi.

\begin{description}
  \item[Diego Ricchiuti] un gioco è un'attività emozionale in cui gli utenti
        interagiscono con una struttura \textbf{ARBITRARIA}, in cui
        \textbf{PERCEPISCONO} un obiettivo e un conflitto. \emph{Arbitrario}
        indica qualcosa di non equo, deciso da qualcuno che ha potere assoluto,
        che può essere il gioco stesso o il designer. \emph{Percepiscono}
        significa che il conflitto può anche non essere presente, ma i
        giocatori sentono che c'è: sentono di farlo per una qualche ragione.
  \item[Mattia Podini] un gioco è un sistema basato su regole che evoca
        \textbf{estetiche} di gioco tramite l'interazione (dove, per ora,
        ``estetiche'' va inteso come sinonimo di emozioni; il termine verrà
        formalizzato più avanti).
  \item[Jack Danna] un gioco è un'attività gestita da regole che coinvolge
        l'utente tramite il raggiungimento di un obiettivo.
\end{description}

Vengono mostrate anche le definizioni di due ex docenti dell'accademia (che
erano stati docenti di game design e level design dello stesso direttore,
quando era studente):
\begin{description}
  \item[Paolo, level designer] un'attività coinvolgente in cui gli utenti
        interagiscono in modo arbitrario con una struttura e percepiscono un
        obiettivo che si risolve in modo ineguale. Come si vede, molto vicina
        a quella del direttore (``è stato convertito'').
  \item[Valerio] uno strumento di \textbf{esplorazione e apprendimento}
        (\emph{a tool of exploration and learning}), dove ``esplorazione'' non
        è fisica ma significa sperimentare, che usa regole, obiettivi e
        simulazioni del mondo reale \textbf{abbassando il rischio}.
\end{description}
Quest'ultima è molto interessante perché coglie un elemento del \emph{play}:
\textbf{abbassare la posta in gioco facilita il gioco}. Chi può ``giocare''
in borsa? Chi è ricco: se perde i soldi, poco importa.

\subsection{La definizione come bussola}

Che cosa farsene, allora, delle definizioni scritte in aula, e perché
dedicarvi un'ora? Per due ragioni. La prima è stata vista: capire che il gioco
non è definibile. La seconda è più personale:
\begin{quote}
\textbf{La vostra definizione di gioco definisce voi come designer.}
\end{quote}
Va riletta, studiata, analizzata:
\begin{itemize}
  \item chi ha scritto \emph{storia} ha un animo narrativo;
  \item chi ha scritto \emph{sistemi} è un futuro \emph{system designer};
  \item chi ha scritto \emph{multiplayer} è interessato a quell'ambito;
  \item chi ha scritto \emph{avventura} è chiaramente interessato ai giochi
        d'avventura;
  \item chi ha scritto \emph{emozioni} vede il lato più romantico ed
        emozionale del gioco;
  \item chi ha scritto \emph{regole} vede la componente più classica del
        mestiere, quella del game designer che sceglie le regole del gioco;
  \item chi ha scritto \emph{personaggi} è interessato alla simulazione delle
        interazioni umane tramite personaggi: non la storia del grande eroe
        che salva il mondo, ma l'umano, il piccolo, le persone che si
        confrontano, le loro vicissitudini;
  \item chi ha parlato di \emph{risolvere problemi} tende, probabilmente, più
        al game design che ad altre discipline.
\end{itemize}
Sono tutte giuste. Sono utili perché diventano una \textbf{bussola}: un punto
di riferimento per i progetti in accademia, per la carriera futura e per il
tipo di designer che si vorrà essere. È inutile che chi ha scritto
``narrativa'' o ``storia'' si metta a fare giochi sistemici: non è ciò che gli
interessa.

La raccomandazione pratica è di \textbf{riscrivere la propria definizione
ogni tanto} (non ogni giorno: ogni anno, ogni biennio o più), \emph{senza
rileggere quella vecchia}. Le cose evolvono, i gusti cambiano, si giocano
giochi diversi, si fanno esperienze diverse, si cresce come designer (che non
significa necessariamente migliorare, ma cambiare).

La definizione del direttore mette in risalto, addirittura in maiuscolo, due
parole estremamente teoriche: \emph{arbitrario} e \emph{percepito}. Non è un
caso: è un designer che adora la \textbf{teoria del game design}, che studia
libri di game design, psicologia, sociologia e antropologia, che fa ricerca
per spingere la teoria un passo più in là. La sua definizione è legata alla
percezione, al voler andare oltre ciò che si è già visto. Questo non la rende
migliore: la rende una bussola per lui. Se tra qualche anno cambierà (e
succederà sicuramente), rappresenterà il suo nuovo modo di affrontare il
problema.

% ---------------------------------------------------------------------
\section{Gli ingredienti della zuppa}
\label{sec:zuppa}

\subsection{La metafora della zuppa}

Non tutto è perduto. Se non si può definire il gioco, se ne possono definire
le \textbf{caratteristiche}. Il gioco è come una zuppa: che cos'è una zuppa?
Acqua con delle cose dentro? Una ricetta? Anche senza sapere definire con
esattezza se qualcosa sia una zuppa, se ne possono \textbf{elencare gli
ingredienti}, e da quelli capire se quella zuppa piace o no, o se è davvero una
zuppa. Lo stesso vale per il gioco: anche se non lo si può definire, se ne
possono definire le \textbf{componenti}, e attraverso la loro analisi capire
se un gioco rientra nei propri gusti, oltre a molte altre caratteristiche che
si scopriranno più avanti.

\subsection{Raph Koster: artefatti ludici e sistemi ludici}

Il primo esempio viene da Raph Koster, autore di \emph{A Theory of Fun for
Game Design} e designer noto soprattutto per \emph{Ultima Online}, uno dei
primi MMO della storia. Koster, invece di definire il gioco, ne distingue le
componenti partendo da due categorie:

\begin{description}
  \item[Artefatti ludici] (\emph{ludic artifacts}): sistemi \textbf{progettati
        per essere giocati} (\emph{designed for play}). Esempi: il \emph{go}, lo
        splendido gioco asiatico; \emph{Magic: The Gathering}; il football e il
        calcio; \emph{Werewolf} (in Italia noto come \emph{Lupus in
        Tabula}); \emph{Space Invaders}; \emph{League of Legends}.
  \item[Sistemi ludici] (sistemi \emph{suscettibili di gioco}, cioè con cui
        si può interagire giocando): sistemi che \emph{non} sono giochi, eppure
        possono essere giocati.
\end{description}

Gli esempi della seconda categoria sono illuminanti:
\begin{itemize}
  \item \textbf{il mercato azionario} (\emph{stock market}): può essere uno
        strumento per fare soldi, o un modo per passare il tempo, se fare soldi
        non è l'obiettivo;
  \item \textbf{uno strumento musicale}: si può suonare Beethoven al violino,
        oppure strimpellare una chitarra, ridere della propria incompetenza e
        far ridere gli altri, con sé e di sé;
  \item \textbf{la guerra}: per chi la vive è tutt'altro che un gioco, ma
        esistono i giochi di guerra, e spesso chi fa la guerra per lavoro (i
        generali) la deve ``deumanizzare'', valutarla come se fosse un gioco,
        perché pensando alle vite umane non si riuscirebbero a fare le scelte
        logicamente corrette;
  \item \textbf{la politica}: il direttore cita un episodio recente di una
        persona eletta consigliere comunale negli Stati Uniti dopo una campagna
        nata come scherzo (``votatemi, non sono competente, ma ne ho voglia'').
        È un caso estremo, ma si dice spesso che i politici ``giocano'' con la
        vita delle persone;
  \item \textbf{l'orecchio interno}: chi gira su se stesso (o lo ha fatto da
        bambino) sta usando la propria conoscenza, empirica o meno, di come
        funziona l'orecchio interno per crearsi un'emozione, la sensazione di
        essere storditi (\emph{dizzy}). Sta ``giocando'' il proprio orecchio
        interno. Il direttore lo fa ancora con i nipotini: vince chi gira di più
        su se stesso e poi resta in piedi (e vincono loro);
  \item \textbf{la fisica}: la palla, l'oggetto più basilare del mondo, ha una
        reazione fisica ed è una forma di gioco con la fisica dell'universo.
        Basta cambiare quanto rimbalza la palla per cambiare completamente il
        gioco: si provi a giocare a calcio con il pallone da basket, a basket
        con il pallone da football americano, a tennis con il pallone da
        calcio. Non si è cambiata nessuna regola, solo l'interazione fisica
        dell'oggetto, eppure è cambiato tutto. Lo stesso vale per le
        palline di gomma dura grandi come una biglia che rimbalzano
        all'infinito (e si perdono ovunque), o per il salto della corda, dove
        la reazione fisica è affidata alle competenze del proprio corpo.
\end{itemize}

\paragraph{Anche i sistemi hanno regole.} Mattia aggiunge che anche questi
sono sistemi, e quindi hanno regole interne. Nella fisica c'è la gravità, ma
c'è anche lo spazio: lo si assume infinito, e ci sono simulazioni che
funzionano assumendolo tale, ma a livello computazionale c'è sempre un limite.
Un esempio che gli studenti sperimenteranno presto: nelle prime lezioni di
level design, quando si impara a modificare il \emph{character}, quasi tutti
aumentano a dismisura velocità di movimento, altezza del salto e \emph{air
control}, ottenendo personaggi che saltano letteralmente \emph{fuori dal
mondo}. Lo spazio dell'engine è finito (una sorta di cupola), e con un salto
abbastanza potente lo si vede rimpicciolirsi mentre si vola ``verso l'infinito e
oltre''. ``Lo farete, lo fanno tutti, è normale.''

\subsection{Game e play sono la stessa cosa}

Perché questa distinzione? Anche Koster non ha una definizione giusta (nessuno
ce l'ha e nessuno mai ce l'avrà), ma la sua divisione mostra che esistono
\textbf{sistemi progettati per il play} (il gioco ``classico'', come lo hanno
definito gli studenti) e una \textbf{mentalità di play} (il gioco inteso in
senso largo, lo stato mentale: tutto può essere un gioco, a patto di essere
nello stato mentale di giocarlo).

Qual è quello giusto? Sono \textbf{la stessa cosa}. Non due facce della stessa
moneta: la stessa cosa. Comprenderlo è \textbf{il primo segreto per sviluppare
giochi}: il gioco è un gioco, ma lo stato mentale è parte del gioco, anzi è il
gioco stesso. C'è sempre uno stato mentale legato al gioco, e \textbf{la
capacità di lavorare con lo stato mentale del giocatore} sarà parte del lavoro
del designer. Più si andrà avanti, nel corso dell'anno e negli anni
successivi, più si parlerà di psicologia, proprio perché lo stato mentale del
giocatore è parte integrante dell'esperienza. Detto in inglese: uno è il
\emph{game}, l'altro è il \emph{play}.

% ---------------------------------------------------------------------
\section{Il framework di Jesse Schell: la tetrade elementare}
\label{sec:schell}

\subsection{Che cosa sono i framework}

Si entra ora ``dentro la zuppa'', con strumenti che da qui in poi si chiameranno
\textbf{framework}: elenchi strutturati delle componenti di un gioco. I primi
due sono quelli di Jesse Schell, autore di \emph{The Art of Game Design: A Book
of Lenses} (il testo di riferimento del corso, si veda la
sezione~\ref{sec:risorse}), e di Brenda Romero e Ian Schreiber, autori di
\emph{Challenges for Game Designers}. Sono tre nomi tra i più importanti della
storia del design: Schell, presentato in lezione come designer di grande
esperienza anche nel gioco da tavolo;\footnote{Jesse Schell è noto soprattutto
per il lavoro ai Disney Imagineering VR Studio, per la fondazione di Schell
Games e per l'insegnamento alla Carnegie Mellon University.} Brenda Romero,
una delle più importanti designer di videogiochi; Ian Schreiber, uno dei più
importanti teorici del design. Le loro versioni del ``contenuto della zuppa''
sono completamente differenti.

Schell divide il gioco in quattro elementi, la cosiddetta \textbf{tetrade
elementare}: meccaniche, narrativa, estetiche e tecnologia.

\subsection{Meccaniche}

Le \textbf{meccaniche} sono le procedure e le regole del gioco. Descrivono
l'\textbf{obiettivo}, come i giocatori possono e non possono raggiungerlo e che
cosa succede quando ci provano. Comprendono:
\begin{description}
  \item[Le regole] ciò che si può e non si può fare.
  \item[I verbi] un altro modo di chiamare le meccaniche: un verbo è
        letteralmente un verbo (saltare, sparare, mirare). Le azioni del gioco
        si esprimono in verbi in \emph{-are}, \emph{-ere}, \emph{-ire}.
  \item[I pezzi] nessuna definizione della classe li aveva citati. I pezzi non
        sono solo il pedone degli scacchi: è un pezzo del gioco anche il
        controller, o il mouse. Gli FPS ne suggeriscono uno, i giochi di guida e
        d'azione un altro, al punto che certi controller vengono venduti a
        centinaia di euro perché danno un vantaggio (il direttore ne mostra uno
        con i tasti posteriori). Il termine corretto, come si vedrà, è
        \emph{tecnologia}.
\end{description}
Schell dice quindi: la prima cosa da definire è l'obiettivo, poi come i
giocatori lo raggiungono, che cosa possono e non possono fare, che cosa succede
quando ci provano, le regole, i verbi e i pezzi.

\subsection{Narrativa}

La \textbf{narrativa} è la sequenza di eventi che si sviluppa nel gioco, che
può essere:
\begin{itemize}
  \item \textbf{lineare e prescritta} (\emph{scripted});
  \item \textbf{ramificata}: non ha un unico finale e, in base a determinate
        scelte del giocatore, può prendere una strada o un'altra
        (\emph{Baldur's Gate 3} è, da questo punto di vista, ``completamente
        fuori di testa'', con un numero enorme di combinazioni di finali);
  \item \textbf{emergente}: a grandi linee, dipende ed è attivamente
        influenzata dalle azioni del giocatore nel corso del gioco. Il concetto
        è molto più ampio e verrà trattato in una lezione specifica.
\end{itemize}

Due osservazioni importanti.
\begin{enumerate}
  \item Schell parla della \emph{sequenza di eventi che si sviluppa nel
        gioco}: non si riferisce solo ai giochi narrativi. Un gioco di auto
        non ha la narrativa percepibile di un \emph{Life is Strange}, ma ha
        comunque una forma di narrativa interna.
  \item Schell usa il termine \textbf{narrativa}, mai \emph{storia}. La
        narrativa implica moltissime cose: la storia, ma anche i personaggi, il
        mondo, tutto ciò che riguarda una descrizione. Alcuni giochi hanno una
        storia, altri un mondo, altri una storia emergente.
\end{enumerate}

\subsection{Un principio fondamentale sui framework}

Da qui discende una regola fondamentale:
\begin{quote}
\textbf{Per essere valido, un framework deve potersi applicare a qualsiasi
gioco.}
\end{quote}
Non può esistere un framework che funzioni solo per un tipo di gioco. Tutto
quello che Schell dice vale tanto per i giochi da tavolo quanto per i
videogiochi. Questo perché Schell si allontana dalla \emph{definizione},
facendo mentalmente lo stesso percorso fatto in aula:
\begin{itemize}
  \item lo scopo di una \textbf{definizione} è \emph{chiudere}: delimitare la
        parola, darle dei limiti;
  \item lo scopo di un \textbf{elenco di elementi} è l'opposto:
        \emph{ingrandire}. Se domani Schell pubblicasse un secondo volume,
        potrebbe aggiungere un quinto elemento, se lo ritenesse necessario.
        \textbf{Non esiste un limite agli ingredienti della zuppa}: la zuppa
        cresce in base alla necessità.
\end{itemize}

Mattia aggiunge una precisazione: è vero che i framework sono fatti per
raccogliere sotto il loro cappello quanti più giochi e metri di analisi
possibili, ma resta vero che spetta alla \textbf{sensibilità del designer}
scegliere, in base a ciò che vuole fare, a ciò in cui è bravo, a dove lavora,
alle risorse e ai limiti (si veda la sezione~\ref{sec:limiti}). È anche il
motivo per cui esistono due framework che dicono più o meno la stessa cosa, con
liste diverse e in parte complementari. Quello di Schell, ad esempio, con la
sua attenzione alla narrativa, si presta meglio a spiegare un gioco
\emph{story driven}; per confrontare due giochi con la stessa tecnologia (ad
esempio giocati sulle stesse periferiche) è più adatto quello di Romero e
Schreiber. Bisogna iniziare a pensare in termini di \textbf{malleabilità} dei
framework: quando conviene usarne uno e quando un altro.

\subsection{Estetiche}

Attenzione: si incontreranno \textbf{due accezioni diverse} del termine
\emph{estetica}. Quella di Schell è la più ampia: le estetiche sono
\textbf{come il gioco si vede, si sente, profuma, si assapora}. Profumo e sapore
non sono un'iperbole: esistono gli scacchi alcolici, in cui ogni pezzo mangiato
è un bicchierino da bere; il \emph{chess boxing}, che alterna partite a scacchi
e round di pugilato; i libri di Geronimo Stilton con le pagine profumate. In
generale: \textbf{come il gioco viene percepito}.

Le estetiche sono un aspetto incredibilmente importante perché sono
\textbf{la relazione più diretta del giocatore con l'esperienza}: il modo più
diretto con cui il giocatore interagisce con il gioco, e la parte ``finale'' di
ciò che percepisce. Da giocatori si vede come il gioco si presenta, quanto è
piacevole giocarlo, quanto è difficile compiere un'azione, che cosa significa
lanciare una palla di fuoco in \emph{Super Mario}: vedere la palla di fuoco.
Quella però è la \textbf{vestizione} applicata sopra le meccaniche. La
meccanica è ``se A, allora succede B''; come B si presenta, di che colori è,
che suono fa, se c'è o no la vibrazione (\emph{rumble}, cioè feedback aptico
sul controller), tutto questo ricade nel regno delle estetiche. Più avanti si
vedrà una terminologia più formalizzata che separa ancora meglio i livelli di
estetica e di meccanica.

\paragraph{Esempio: il coltello.} Si immagini il proprio gioco preferito, ma
ogni volta che si subisce un danno qualcuno compare dal nulla e accoltella il
giocatore. Probabilmente non sarebbe più il gioco preferito. È un caso
ridicolo, ma mostra che se la relazione estetica del gioco fosse il dolore,
quasi tutti abbandonerebbero (non tutti, però, perché c'è chi apprezza il
dolore). Gli elementi visivi, interattivi, perfino fisici sono parte
dell'esperienza.

\paragraph{Esempio: le miniature.} Chi gioca con le miniature si è punto
innumerevoli volte con le lance (ad esempio di \emph{Warhammer Fantasy}),
soprattutto in fase di montaggio: un misto di amore e odio che fa parte del
percorso. Una miniatura di metallo, di plastica, di plastica con un leggero
strato metallico magnetizzabile, di ghisa: ognuna dà sensazioni diverse.
Lo stesso vale per le carte \emph{foil}, con la loro sensazione al tatto e la
loro piegatura. Sono tutte \textbf{sensazioni tattili}, elementi fondamentali.

\paragraph{Esempio: Raggio Aurora.} In \emph{Pokémon} molte mosse sono, dal
punto di vista meccanico, semplicemente una statistica di danno e un tipo:
Braciere (fuoco) e Pistolacqua (acqua) sono meccanicamente quasi identiche,
cambia solo il tipo. Eppure tutti hanno mosse preferite, non perché più
efficaci, ma perché l'animazione è spettacolare. Il direttore, che ha giocato
\emph{Pokémon Rosso e Blu}, ricorda ancora \textbf{Raggio Aurora}, la cui
animazione (``il raggio epilessia'') è diventata sempre più vistosa di
generazione in generazione. Raggio Aurora è una mossa mediocre: in competitivo
nessuno la usa, perché esiste Geloraggio, che ne è la versione migliore.
Eppure è la mossa che ricorda, per la sua \textbf{componente estetica}, che per
i suoi gusti (non in assoluto) la rende più interessante di altre.

\subsection{Tecnologia}

La \textbf{tecnologia} è tutto il materiale, tecnologico o fisico, che riguarda
il gioco: il computer su cui gira, l'engine che lo fa girare, il controller, il
materiale delle miniature. È importante perché cambia l'esperienza:
\begin{itemize}
  \item il \emph{feeling} di una carta o di una miniatura cambia da materiale a
        materiale; giocare con un controller o con mouse e tastiera è diverso;
  \item se si fa girare \emph{GTA V} su un vecchio portatile a 6 fps,
        l'esperienza è pessima: del gioco non è cambiato nulla, tranne il fatto
        che gira a 6 fps. È la ragione dei \textbf{requisiti minimi} e dei
        \textbf{requisiti consigliati} per un'esperienza decente;
  \item in \emph{Pokémon GO}, che usa la posizione reale del giocatore, alcuni
        streamer sono stati derubati del telefono dopo essersi avventurati, per
        catturare un Pokémon raro, in zone poco sicure delle città: casi
        estremi, ma che mostrano quanto la tecnologia influenzi il gioco;
  \item in \emph{Invizimals} (su PSP) si acquistavano in edicola delle carte o
        miniature dei mostri, che venivano poi ``catturati'' inquadrandoli con
        la fotocamera della console: anche questa è tecnologia.
\end{itemize}

\paragraph{Esempio: Commander e le cento carte.} \emph{Magic} nasce con mazzi
da 60 carte. \emph{Commander} è un formato creato dalla comunità, poi diventato
famoso e consolidato, che si gioca con mazzi da \textbf{100 carte}. Qui emerge
un limite della tecnologia di Magic: \textbf{mescolare cento carte è
estremamente difficile e noioso}. In un gioco in cui mescolare è un'azione
abbastanza frequente, questo diventa parte di un'esperienza negativa, tanto che
esistono tavoli in cui i giocatori si accordano per non inserire carte che
fanno mescolare il mazzo, così da farlo solo all'inizio. Il problema nasce dal
fatto che Magic è stato pensato con carte di un certo peso, dimensione e costo,
che funzionano in mazzi da 60, non da 100. Si potrebbe sostenere che non sia un
errore di Wizards of the Coast, dato che il formato non è nato da loro; ma il
punto è che la tecnologia cambia l'esperienza.

% ---------------------------------------------------------------------
\section{Il framework di Brenda Romero e Ian Schreiber}
\label{sec:romero-schreiber}

\subsection{Perché più punti di vista}

Si guarda ora il contenuto della zuppa dal punto di vista di Romero e
Schreiber. Non perché siano più bravi di Schell, o viceversa, né perché la
visione del direttore sia migliore delle loro: \textbf{non esiste un framework
migliore}, così come la divisione di Koster non è migliore o peggiore delle
altre. Mostrare la stessa cosa da più punti di vista serve a capirne la
complessità; inoltre ognuno tenderà per natura a sentirsi più affine a uno o
all'altro, perché ognuno ha una propria visione creativa, e lavorare con uno
strumento allineato al proprio modo di pensare è più facile. Conoscere modi di
pensare diversi serve proprio a questo.

\subsection{Gli elementi}

\begin{description}
  \item[Game qualities] (qualità di gioco) al posto delle meccaniche, Romero
        definisce i \textbf{sistemi di gioco}. Invece di concentrarsi su
        saltare, mirare, sparare, ricaricare, schivare (\emph{dodge}), parare,
        attaccare (tutte meccaniche, tutti verbi), si ragiona per gruppi:
        ad esempio la qualità \emph{shooting} include tutte le caratteristiche
        legate alle meccaniche di tiro, che non necessariamente sono tutte
        presenti in un dato gioco (in \emph{Doom} non si ricarica, ma è
        comunque uno shooter). Altri esempi sono \emph{PvP} e \emph{PvE}: se un
        gioco è PvP ha il multiplayer, e non serve definirne la meccanica, è
        sottinteso; se ha il PvP avrà probabilmente un sistema di barre della
        vita o qualcosa di simile, e così via.
  \item[Giocatore] chi gioca? Perché gioca? Che attività sta svolgendo? Che
        cosa desidera, che caratteristiche ha? È un elemento fondamentale, su
        cui si tornerà in altre lezioni: \textbf{il giocatore è parte
        integrante del gioco}. Se si mette il direttore a giocare a
        \emph{Canasta}, non gli piace: non è il suo gioco. Questo non significa
        che Canasta sia un brutto gioco, ma che non è il gioco per lui, pur
        avendo la stessa tecnologia (le carte) di giochi che ama e molte
        caratteristiche simili (interazione con altri giocatori, combinazioni
        di carte, strategia). Allo stesso modo, nessuno gioca a \emph{Tekken}
        ``sul serio'' contro un nipotino di tre o quattro anni: si sceglie
        piuttosto un gioco di macchinine, o \emph{Mario Kart} rendendogli le
        cose più facili. La \textbf{difficoltà} dipende dal giocatore.
  \item[Tema] tutti gli elementi legati alla storia e/o al
        \textbf{significato} del gioco. È simile alla narrativa di Schell, ma
        aggiunge il significato, che include ad esempio la comprensione di un
        sottotema filosofico o narrativo del prodotto.
  \item[Obiettivo] lo scopo del gioco: che cosa si cerca di ottenere e di
        raggiungere. Molti studenti lo avevano inserito nella propria
        definizione.
  \item[Regole] qui sono una struttura separata, che definisce soprattutto
        \textbf{ciò che il giocatore non può fare}. Nelle meccaniche di Schell
        le regole definivano anche ciò che si può fare, ciò che si vuole
        ottenere e che cosa succede; qui ciò che si può fare è già definito
        nelle \emph{game qualities}.
  \item[Risorse] (illustrate da Mattia) sono i \emph{game pieces} che possono
        essere manipolati, usati, consumati, collezionati o modificati:
        elementi che possono cambiare, progredire o essere spesi all'interno
        dell'esperienza. In \emph{The Witcher} sono risorse i soldi, ma anche
        armature e armi, e anche le carte di \emph{Gwent}, con cui si
        costruisce, migliora e potenzia un mazzo.
  \item[Interazione] la \textbf{sequenza di causa ed effetto}.
\end{description}

\subsection{Confronto con Schell}

Molto di ciò che Schell raccoglie in un unico punto, le meccaniche (con
regole, verbi e pezzi), qui è sparso in più elementi: le \emph{game qualities}
lavorano a un livello più alto, le regole a un livello più basso (con
particolare attenzione a ciò che non si può fare); regole e risorse insieme
iniziano a dare un quadro molto simile alle meccaniche di Schell.

C'è però un rischio: nei giochi moderni, l'elenco delle risorse può essere
\textbf{potenzialmente infinito}. In un \emph{Final Fantasy} ogni pozione è
una risorsa, ogni attacco, ogni personaggio, ogni punto vita, ogni punto
esperienza: è facile perdersi in una lista senza fine. Così, mentre sulle
\emph{game qualities} Romero sembra fare un lavoro di ``asciugatura'', sulle
risorse il framework esplode e vuole definire ogni singolo elemento.

Anche l'\textbf{interazione} è interessante: in Schell causa ed effetto
restano impliciti (se si definisce che cosa si può e non si può fare, come lo
si fa e che cosa si vuole ottenere, si sta facendo come nel calcio, lasciando
implicito il contropiede); qui invece si vogliono definire: ``il giocatore
potrebbe fare il contropiede''.

\paragraph{Esempio: il camping.} Nei giochi multiplayer il \emph{camping} è
l'atto di appostarsi in un punto particolarmente vantaggioso e strategico,
allo scopo di eliminare gli altri giocatori che passeranno da lì, riducendo al
minimo il proprio rischio. È una strategia validissima, solitamente
considerata negativa perché ritenuta meno \emph{skillful} (non è vero, ma è
così che viene percepita). Nel framework di Romero e Schreiber il camping va
inserito tra le interazioni; in Schell resta non definito. Il che non
significa che non esista: semplicemente non è esplicitato.

Ancora una volta: nessuno dei due è migliore o peggiore. Sono due punti di
vista.

% ---------------------------------------------------------------------
\section{Il framework di Stéphane Bura}
\label{sec:bura}

\subsection{Una tabella a incroci}

Il terzo framework, introdotto da Jack, è quello di Stéphane Bura. A
differenza dei precedenti, \textbf{si legge come una tabella a incroci}: non
basta considerare i singoli termini, bisogna leggerli in combinazione. Sulle
righe ci sono quattro \textbf{livelli}, sulle colonne tre \textbf{dimensioni}.

\begin{description}
  \item[Action] (azione) ciò che il giocatore fa, ciò che ha la possibilità di
        compiere. È il livello più basso, quello del \emph{second to second},
        del gameplay più rapido: la chiave che apre una porta, l'attacco base.
  \item[System] (sistema) ciò che fa il gioco. È un livello superiore: finisco
        questo bioma e passo a un altro, vado di qua o di là.
  \item[Self] (sé) l'interazione del giocatore come singolo individuo con il
        gioco: interazioni che non si limitano a sfruttare il sistema (quello è
        \emph{action}), ma ne fanno un \textbf{uso creativo o strategico}.
  \item[Social] (sociale) tutto ciò che è legato alla comunità: il giocatore
        immerso in una community, la sua capacità di relazionarsi con gli
        altri.
\end{description}
Ciascuno di questi livelli si articola in tre dimensioni: \textbf{Freedom}
(libertà), \textbf{Mastery} (maestria, padronanza) e \textbf{Data} (dati,
informazioni).

\begin{table}[htbp]
  \centering
  \small
  \begin{tabular}{@{}p{0.11\textwidth}p{0.27\textwidth}p{0.27\textwidth}p{0.27\textwidth}@{}}
    \toprule
    & \textbf{Freedom} & \textbf{Mastery} & \textbf{Data} \\
    \midrule
    \textbf{Action} & opportunità, abilità
                    & competenza nelle azioni a basso livello, spesso fisica
                    & risorse del mondo di gioco, collezionabili, regole
                      operative \\
    \addlinespace
    \textbf{System} & esplorazione, sperimentazione, scopo
                    & competenza sui sistemi complessi
                    & regole e loro comprensione, sistemi che trasformano le
                      risorse \\
    \addlinespace
    \textbf{Self}   & strategia, creatività
                    & etica, sensibilità personale nell'uso del sistema
                    & \emph{mementos}, achievement, memoria \\
    \addlinespace
    \textbf{Social} & community, supporto, esperienza condivisa
                    & abilità sociali (mentire, leggere, collaborare)
                    & meta, metagaming \\
    \bottomrule
  \end{tabular}
  \caption{Il framework di Stéphane Bura come illustrato in lezione.}
  \label{tab:bura}
\end{table}

\subsection{La colonna Freedom}

\begin{description}
  \item[Action $\times$ Freedom: opportunità e abilità] una delle
        caratteristiche tipiche di un gioco è dare al giocatore delle abilità e
        lasciarlo libero di sceglierle. Il ``mega pugno della morte'' si attiva
        con L2 più quadrato quando si vuole: si è liberi, ed è un'azione.
  \item[System $\times$ Freedom: esplorazione, sperimentazione, scopo] si è
        liberi di andare dove si vuole, ma \emph{entro i limiti del gioco}. Se
        c'è un muro invalicabile, il muro è invalicabile; se il personaggio non
        sa nuotare, in acqua non si va. Il sistema limita la libertà.
  \item[Self $\times$ Freedom: strategia e creatività] in \emph{Minecraft} il
        \emph{mine} (dare un pugno e rompere qualcosa) è un'azione. Capire che
        scavando un fossato abbastanza largo e profondo i creeper ci cadono
        dentro di notte, e quindi non si muore, non è un'abilità (non esiste
        un tasto che fa comparire un fossato): è un'\textbf{applicazione
        strategica o creativa del sistema}. Lo stesso vale per l'uso creativo
        di Minecraft: c'è chi ha ricostruito in scala 1:1 il continente di
        Westeros del \emph{Trono di Spade}, cosa che nulla nel gioco chiede di
        fare. Jack cita anche \emph{The Uncensored Library}, un progetto di
        Reporters sans frontières: una biblioteca costruita in Minecraft che
        raccoglie articoli e testi censurati in vari paesi, consultabili
        proprio perché si trovano dentro un gioco.
  \item[Social $\times$ Freedom: community, supporto, esperienza condivisa]
        la libertà di entrare in una comunità, interna al gioco (una gilda) o
        esterna (un canale Discord), e di interagire attraverso di essa.
\end{description}

\subsection{La colonna Mastery}

\begin{description}
  \item[Action $\times$ Mastery] tutto ciò che riguarda lo sviluppo di
        competenza nelle azioni a basso livello, che a volte sono abilità molto
        fisiche: saper eseguire un attacco, combinare dei tasti. I
        \emph{rhythm game}, come \emph{osu!} o \emph{Dance Dance Revolution},
        richiedono proprio questo tipo di maestria fisica.
  \item[System $\times$ Mastery] la competenza sui sistemi complessi. In
        \emph{Pokémon} è la capacità di fare \emph{team building}: costruire
        una squadra bilanciata con un attaccante fisico, un attaccante speciale,
        un ``muro''. In un gioco di carte come \emph{Yu-Gi-Oh!} è sapere la
        ripartizione classica del mazzo (ad esempio 20 mostri, 10 magie, 10
        trappole). Quanto bene si sono compresi i sistemi e come li si può
        sfruttare.
  \item[Self $\times$ Mastery] un discorso legato a una certa
        \textbf{etica}, cioè alla sensibilità personale nell'uso del sistema
        (si veda l'esempio di Shudderwock qui sotto). Vi rientrano anche le
        creazioni complesse, come quelle di Minecraft.
  \item[Social $\times$ Mastery] quanto si è bravi nelle relazioni sociali e
        nello sfruttare abilità sociali. In \emph{Among Us} e in qualsiasi
        gioco di deduzione sociale: quanto si è bravi a mentire, a far credere
        di non avere un certo ruolo o di essere amichevoli; ma anche a leggere
        gli altri e a collaborare con loro.
\end{description}

\paragraph{Esempio: Shudderwock in Hearthstone.} In \emph{Hearthstone} le
carte si giocano in un ambiente digitale, con animazioni, e ogni turno ha un
limite di tempo. Quando fu introdotta la carta \textbf{Shudderwock}, il suo
effetto era: quando viene giocata, riattiva tutti i \emph{grido di battaglia}
(\emph{battlecry}, gli effetti che si attivano quando una carta viene messa in
campo) giocati fino a quel momento nella partita. Nacque così un mazzo
specifico, non previsto, riempito di carte con grido di battaglia; spesso lo
Shudderwock veniva anche riattivato una seconda volta, facendo scattare tutti
quegli effetti due volte. Le animazioni erano così lunghe da sforare nel turno
dell'avversario, ``mangiandoselo'': chi giocava il mazzo poteva finire il
proprio turno, mentre le animazioni continuavano e l'avversario di fatto non
giocava. Una volta conosciuta questa cosa, \textbf{quanto è etico giocare quel
mazzo}, sapendo che è estremamente facile vincere sfruttando una lacuna, un
bug o qualcosa che gli sviluppatori non avevano previsto e che probabilmente
verrà corretto? Questa è la dimensione \emph{Mastery Self}. (Il problema fu
poi risolto: Shudderwock fu indebolito e il sistema del tempo modificato.)

Jack aggiunge che una recente espansione di Hearthstone ha introdotto una
meccanica, chiamata in italiano \emph{Innesto}, che divide l'effetto dei
servitori in due parti, la \emph{condizione di attivazione} e la
\emph{conseguenza}, e permette di fonderne più di una. Ne nascono combinazioni
estreme: il classico Doomsayer (``all'inizio del tuo prossimo turno distruggi
tutti i servitori'') fuso con un grido di battaglia diventa ``grido di
battaglia: distruggi tutti i servitori''; oppure ``ogni volta che peschi una
carta, pesca una carta'', con effetti facili da immaginare (un ``Vaso
dell'Avidità'' potenziato all'infinito, per citare \emph{Yu-Gi-Oh!}).

\paragraph{La difficoltà nei framework.} Uno o due studenti avevano citato la
difficoltà, la sfida, nelle proprie definizioni. Non compare esplicitamente né
in Schell né in Romero e Schreiber, mentre in Bura c'è, sotto forma di
\emph{mastery}, cioè la possibilità di diventare più bravi. Significa che per
Schell e Romero la difficoltà non conta? No: la fanno rientrare in altri punti
(nelle meccaniche, nel giocatore). \textbf{Gli ingredienti della zuppa sono
sempre gli stessi}; i designer si concentrano su elementi diversi in base ai
propri interessi, al proprio modo di pensare e di creare.

\subsection{La colonna Data}

\begin{description}
  \item[Action $\times$ Data: risorse, collezionabili, regole operative] le
        monete di \emph{Super Mario}: si raccoglie una risorsa con una
        micro-azione che produce un cambiamento in quelle che Romero chiama
        \emph{risorse}.
  \item[System $\times$ Data] qualcosa di più complesso: le regole e la loro
        comprensione (tutorial, spiegazioni), ma anche i sistemi complessi che
        incidono in modo importante sulle risorse. Giochi come
        \emph{Civilization} si basano interamente sull'acquisizione,
        sull'accumulo e sull'``esplosione'' (aumento esponenziale) delle
        risorse; così molti strategici in tempo reale, come \emph{Age of
        Empires}, basati sulla raccolta di risorse.
  \item[Self $\times$ Data: mementos, achievement, memoria] ciò che il
        giocatore si porta via e che riguarda lui. Spesso sono
        \emph{mementos}, ricordi: la \emph{photo mode}, la propria
        \emph{build}, il personaggio creato in un certo modo nell'editor
        iniziale. Il direttore racconta di aver creato, in un soulslike prima
        di una patch, un personaggio con un naso così lungo da deformare tutti
        gli elmi (che si adattano alle ossa del volto): nessuna funzione,
        nessun effetto sul gioco, ma ``mi sentivo figo''. Ancora più
        chiaramente rientrano qui gli \textbf{achievement}: per chi
        ``platina'' i giochi, ottenerli è parte integrante dell'esperienza.
  \item[Social $\times$ Data: il meta] i dati sociali sono legati al
        \textbf{metagaming}: non gioco quella cosa perché non vince in
        classifica (\emph{ladder}); gioco quel personaggio perché vince; uso
        quel mouse con la rotellina magnetica perché dà certi vantaggi.
\end{description}

\paragraph{Mastery o meta?} Una domanda dall'aula: lo sfruttamento del meta
(\emph{Social Data}) non si confonde con la padronanza dei sistemi
(\emph{Mastery System})? La risposta: la \textbf{mastery richiede una
competenza}, il \textbf{meta richiede una conoscenza}. Se si mette qualcuno al
volante di una Formula 1 Ferrari, la sua capacità di guidarla è la sua
mastery: può crescere nel tempo, e dopo due anni sarà un pilota migliore. Il
fatto che la scuderia gli dica ``questa curva va presa a 197,5 km/h in
quinta'' è meta.

Un altro esempio: quanti costruiscono da soli la propria build o il proprio
mazzo, invece di cercare online il più efficace? È normale divertirsi a
costruire qualcosa di assurdo, ma quando si vuole qualcosa di efficace si
cercano dati online, perché qualcuno ha già fatto i calcoli. Imparare poi a
\emph{usare} quel mazzo è un discorso separato.

Nei picchiaduro (\emph{Mortal Kombat}, \emph{Dragon Ball FighterZ},
\emph{Dragon Ball: Sparking! Zero}, il seguito dei \emph{Budokai Tenkaichi})
il meta può dire, per esempio, che Vegeta Super Saiyan 2 è il miglior
\emph{rusher}, cioè il personaggio migliore per andare addosso all'avversario,
sulla base di dati calcolati (velocità delle animazioni e così via). Ma il
direttore, che in questi giochi usa gli \emph{zoner} (personaggi che combattono
a distanza), perderebbe comunque tutte le partite con quel Vegeta. Conoscere
l'informazione è un vantaggio, perché rende statisticamente più probabile
vincere con qualcosa di matematicamente migliore; ma \textbf{sapere
un'informazione non significa saperla pilotare}.

% ---------------------------------------------------------------------
\section{Il framework del direttore}
\label{sec:framework-direttore}

Per mostrare ancora una volta che non esiste un framework giusto, il direttore
mostra il proprio, che verrà spiegato più avanti nel corso. I suoi elementi
sono: \textbf{player}, \textbf{meccaniche}, \textbf{dinamiche},
\textbf{tema} e \textbf{tecnologia}.

Ogni studente, col tempo, svilupperà il proprio modo di interpretare gli
ingredienti della zuppa: un mix di questi framework e di altre cose, senza
limitarsi a pensare ``non lo aveva scritto Schell, non lo aveva scritto Romero,
quindi non posso inserirlo''. Prima, però, bisogna imparare perfettamente ad
analizzare un gioco con i framework esistenti, perché a questo punto della
carriera, statisticamente, nessuno è ancora Jesse Schell, Brenda Romero o Raph
Koster. All'inizio ci si siede \textbf{sulle spalle dei giganti}
(\emph{standing on the shoulders of giants}): loro hanno già fatto questo
lavoro, e non serve ripartire da zero. Col tempo ognuno avrà il proprio, che
potrebbe rivoluzionare l'industria o essere semplicemente uno strumento
personale di lavoro, come quello del direttore.

% ---------------------------------------------------------------------
\section{Le lenti: non esiste un metodo unico}
\label{sec:lenti}

\subsection{Il principio più importante}

Qual è, dunque, il framework migliore? Nessuno: sono \textbf{lenti}
(\emph{lenses}), ed è il motivo per cui il libro di Schell ha come sottotitolo
\emph{A Book of Lenses}. È il primo e più importante principio da ricordare,
da ora fino alla fine della propria carriera di designer, e non solo di
videogiochi:
\begin{quote}
\centering
\textbf{Non esiste una soluzione. Non esiste un santo Graal. Non esiste un
unico metodo che risolva tutti i problemi.}
\end{quote}
Questo è al tempo stesso \emph{orribile} e \emph{meraviglioso}:
\begin{itemize}
  \item orribile, perché a volte ci si scontra con problemi che sembrano
        insormontabili, e trovarsi davanti a qualcosa di apparentemente
        impossibile non è piacevole;
  \item meraviglioso, perché \textbf{tutto è superabile}: bisogna solo
        cambiare la lente con cui lo si guarda.
\end{itemize}

\subsection{La migliore difesa contro l'intelligenza artificiale}

C'è anche un secondo aspetto: questa è \textbf{la più grande protezione che il
designer ha contro le intelligenze artificiali}. Perché un'AI è brava a
programmare? Perché molti problemi di programmazione si risolvono in un modo
noto. L'ordinamento (\emph{sorting}), cioè prendere qualcosa di disordinato e
ordinarlo, ha algoritmi già risolti: esistono e sono quelli. Non esiste invece
un ``sorting delle meccaniche'': non si riesce nemmeno a definire che cos'è un
gioco, figuriamoci a ordinarlo. Se non esiste un'unica soluzione, non esiste un
algoritmo; e allora la creatività e la capacità di trovare metodi diversi
restano il sistema migliore per fare design.

Il rovescio della medaglia è che bisogna imparare \textbf{più metodi
possibili}. Si può avere un metodo preferito, ad esempio Romero, che risolve
il 90\% dei problemi; poi arriverà il problema che non risolve. Allora si prova
Schell; se non funziona, un terzo, un quarto, un quinto, un sesto. I docenti ne
forniranno moltissimi (non si chiede agli studenti di inventarli), e prima o
poi uno funziona. Il prezzo è che, per sfortuna, si può doverli provare tutti e
scoprire che era l'ultimo: ci si sarà messo più tempo del previsto. Per questo
bisogna \textbf{sempre guardare tutto da più punti di vista}.

\subsection{The right tool for the right job}

Come dice Mattia, \emph{the right tool for the right job}: alcuni framework
funzionano meglio con certi giochi, altri con altri. E c'è un ulteriore
discriminante: \textbf{quelli che funzionano meglio con voi}, con cui ci si
trova più a proprio agio. Jack racconta di conoscere molti framework ma di
usarne prevalentemente tre: se nessuno dei tre funziona passa ai successivi,
ma non avrebbe senso non partire da quelli, perché sono diventati il suo
\emph{bread and butter} (``pane e burro'', o ``pane e marmellata''
all'italiana): qualcosa di talmente rapido e naturale da non costare fatica.

\subsection{Il Deck of Lenses}

Il libro di Schell ha anche una versione in forma di mazzo di carte, il
\textbf{Deck of Lenses} (ne esiste anche una versione online). Il
funzionamento è semplice: davanti a un problema, si estrae una carta a caso e
si prova ad analizzare il problema secondo quella lente.

\paragraph{Esempio: la lente della cooperazione.} Si immagini di essere gli
sviluppatori di \emph{Hearthstone} alle prese con il problema di
Shudderwock. Si mescola il mazzo e si estrae la \emph{lente della
cooperazione}. Ci si chiede, ad esempio, se nel problema ci sia abbastanza
\emph{comunicazione}. Ne nasce un'idea: un sistema per cui, se si gioca un
mazzo con Shudderwock, gli altri lo sanno e possono decidere di non
partecipare alla partita conoscendone il rischio. È una buona soluzione? Non
si sa; ma è emersa ponendosi una domanda attraverso una lente.

\paragraph{Esempio: la lente della sfida.} Si estrae un'altra carta, la
\emph{lente della sfida} (\emph{challenge}): ``la sfida è al centro di quasi
tutto il gameplay; si potrebbe perfino dire che un gioco è definito dai suoi
obiettivi e dalle sue sfide''. La domanda da porsi è: \emph{quali sono le
sfide del mio gioco?} La sfida di Hearthstone è costruire un mazzo e
pilotarlo. Il problema di Shudderwock, allora, è che è un mazzo \emph{meta}:
tutti possono copiarlo e applicarlo, togliendo parte della sfida. Una
soluzione (inventata sul momento) potrebbe essere oscurare i siti che
diffondono il meta. Probabilmente è un'idiozia, perché non si possono bloccare
tutti i siti; ma, di nuovo, è emersa analizzando il problema da un punto di
vista specifico.

Questo è lo scopo dei framework forniti nel corso: \textbf{analizzare i
problemi da punti di vista diversi}, scoperti da designer non necessariamente
più bravi degli studenti, ma con più esperienza. Magari un giorno qualcuno
sarà più bravo di Jesse Schell (e ce lo si augura sempre: la qualità del
design mondiale può solo beneficiarne); ma per natura ora si ha meno
esperienza, e sedersi sulle spalle dei giganti fornisce un'intera serie di
lenti.

% ---------------------------------------------------------------------
\section{Il vocabolario del designer}
\label{sec:vocabolario}

\subsection{La metafora del medico}

Mattia introduce l'ultima parte con una metafora. Quando si va dal medico e si
dice ``sto male, mi fa male la pancia, ho un fastidio generico'', il medico non
può aiutare. Il medico \emph{indaga}: tocca, palpa, fa altre domande,
raccoglie dati, cerca di isolare il problema. Nel game design funziona allo
stesso modo. L'obiettivo finale riguarda il \textbf{responso emotivo} del
giocatore: bisogna capire se è quello desiderato, se il giocatore sta provando
le emozioni che si vuole provi in quel momento (e anche che tipo di giocatore
sia). Se la gran parte dei giocatori prova le emozioni desiderate, è un grande
successo; se no, bisogna indagare.

Ma già a monte bisogna stabilire, con il team, \textbf{che cosa il gioco deve
far scaturire nel giocatore}. Dire ``il nostro gioco deve essere divertente''
non basta: non significa nulla, non dà alcuna informazione. Bisogna delimitare
il problema, trovare \textbf{l'emozione specifica} o l'insieme di emozioni
specifiche, di diversa natura e diversa intensità, che il gioco deve far
provare.

\subsection{La parola proibita}

Da qui una regola del corso, tra le più importanti in assoluto: \textbf{non si
usano più le parole ``divertente'', ``intrattenente'' e simili}, né le loro
varianti o sinonimi (\emph{fun}, \emph{immersivo}\dots). Sono parole che non
vogliono dire nulla. La regola vale in qualunque task e in qualunque
conversazione, dentro e fuori dalla lezione. Chi la infrange viene ``punito''
(non fisicamente, ovviamente) con \textbf{task aggiuntivi}, con modalità
diverse a seconda del docente:
\begin{itemize}
  \item Jack applica una punizione \emph{collettiva}: se una persona usa la
        parola proibita, \textbf{tutta la classe} riceve un task aggiuntivo. Il
        motivo è pratico: con gli studenti dell'anno precedente, che facevano
        molta fatica a smettere, ha funzionato proprio perché il peso non
        ricadeva solo sul singolo, ma su tutti;
  \item Mattia punisce il singolo, ma in modo \emph{progressivo}: più volte si
        usa la parola, peggiore è il gioco da analizzare nel task aggiuntivo,
        partendo da un gioco ``mediamente terribile'' fino ad arrivare a titoli
        come \emph{Hatoful Boyfriend} (un \emph{dating sim} con i piccioni) o
        a simulatori di appuntamenti con gli insetti.
\end{itemize}
In aula uno studente racconta che uno studente del terzo anno gli aveva già
consigliato, il primo giorno, di ``dimenticare la parola divertente''; e
nel raccontarlo la pronuncia, guadagnandosi il primo task aggiuntivo del
corso.

Perché tanta severità? Perché la parola proibita è un concetto
\textbf{completamente inutile}: non esprime \emph{perché} un gioco funziona.
Un gioco funziona perché suscita emozioni, e quelle emozioni vanno definite
attraverso un sistema di design. Anzi, come precisa Mattia, è inutile
\emph{nella migliore delle ipotesi}: quando la parola arriva da altre
discipline o, peggio, da altri designer, diventa ancora più difficile capire
qual è il fulcro del problema o qual è l'equivalente emotivo corretto. Mettere
nella \emph{visione} del progetto (concetto che verrà approfondito), tra i
pilastri del gioco, ``deve essere divertente'' non significa nulla e crea
problemi a tutto il team. In più, toglie \textbf{credibilità} a chi parla.

Non si tratta di una stranezza dei docenti: \textbf{parlare come un designer
significa essere riconosciuti come designer}. Il vocabolario è parte integrante
del mestiere, perché il designer è un \textbf{comunicatore}: comunica con il
team, con il giocatore e con altri designer. Si vedranno tecniche di
comunicazione e si imparerà a fare documentazione, che non è altro che una
forma di comunicazione scritta (esistono la comunicazione orale, verbale, non
verbale e scritta).

\subsection{Non solo ``divertente'': anche ``bello'', ``valido'', ``noioso''}

Una studentessa chiede se la parola \emph{bello} sia equivalente. Sì: dire ``la
narrativa è bella'' è altrettanto inutile. \emph{Schindler's List} è un libro
meraviglioso, ma non è ``bello'' nello stesso modo in cui è bello \emph{Harry
Potter}: il primo racconta di Oskar Schindler e del salvataggio di centinaia di
ebrei dallo sterminio nazista; il secondo è apprezzato per la fantasia, per le
sensazioni che suscita. Leopardi scrive in un modo in cui scrive solo Leopardi:
può piacere o non piacere, ma non lo si può definire ``bello'' come il
\emph{Decameron}, perché sono cose completamente diverse. ``Bello'' è
semplicemente un \textbf{giudizio positivo}, e il problema è che non dice
\emph{perché} qualcosa è piaciuto:
\begin{itemize}
  \item di Leopardi è piaciuta la scrittura, o i temi malinconici?
  \item del \emph{Decameron} è piaciuto lo stile comico, o il fatto che
        rompesse le convenzioni sulla sessualità del suo tempo?
  \item di \emph{Schindler's List} è piaciuto il coraggio storico di Schindler
        nell'affrontare i nazisti, o il modo in cui è scritto?
\end{itemize}

Neanche \emph{valido} va bene, né \emph{noioso} (vale esattamente lo stesso al
contrario), né \emph{schifoso}: nessuna di queste parole descrive una
caratteristica. Bisogna esprimere \textbf{che cosa il gioco cerca di ottenere,
se lo ottiene o no, e con quali caratteristiche}. Che una cosa sia buona o
cattiva è irrilevante.

\paragraph{Una tecnica pratica.} Quando viene da scrivere ``questo gioco è
[parola proibita]'', si completi la frase: ``questo gioco è [parola
proibita] \emph{perché} è comico''. Poi si tagli la parte centrale: ``questo
gioco è comico''. Bisogna imparare a parlare \textbf{dei dettagli}, del
\emph{perché} qualcosa funziona, non del \emph{se} è buono o cattivo.

\paragraph{Esempio: Dan Brown.} Al direttore piacciono le storie di Dan Brown,
ma ne detesta profondamente lo stile di scrittura: per questo non ne legge i
libri, ma ne guarda i film. Non può quindi dire che \emph{Inferno} sia
``bello'' o ``brutto'': il libro ha una buona storia ma uno stile di scrittura
che non apprezza; il film andrebbe giudicato su altre caratteristiche, come la
recitazione. Le opere hanno una serie di caratteristiche e di obiettivi che
dipendono dal medium: per i videogiochi la tecnologia, gli fps, la narrativa,
le meccaniche e quindi il gameplay, le qualità di gioco, i giocatori. C'è chi
potrebbe dire di \emph{League of Legends} ``adoro tutto tranne i giocatori''
(cosa comune, visto che la community è spesso considerata tossica); lo stesso
si sente dire di Commander, ``formato splendido, giocatori orribili''.

Mattia chiude il punto ricordando che il gioco non si colloca quasi mai ``nel
vuoto'': va legato ai giocatori e alle loro necessità. Al direttore piacciono
le storie di Dan Brown ma non il suo stile, per come è lui e per le sue
necessità di lettore; è un discorso che si allargherà ad altri argomenti.

\subsection{Il nuovo vocabolario: una prima lista di emozioni}

Al posto della parola proibita si comincia a usare un vocabolario di emozioni
specifiche. La lista seguente \textbf{non è esaustiva} (ne esistono molte di
più, e verrà ampliata), ma è un punto di partenza.

\begin{description}
  \item[Excitement (eccitazione)] affrontare pericoli, andare veloci, effetti
        luminosi (il ``raggio epilessia'' di Raggio Aurora), suoni forti: tutto
        ciò che fa pompare il sangue nelle vene e crea uno stato di
        eccitazione.
  \item[Challenge (sfida)] Malenia: è difficile, voglio batterla, devo
        riuscirci. Una sfida può essere anche mentale, non solo fisica.
  \item[Competition (competizione)] io contro Mattia: chi fa di più, chi
        vince. Molti giochi, soprattutto non digitali (che di solito prevedono
        almeno due giocatori), si basano sulla competizione. \textbf{Attenzione
        alla differenza}, che negli anni ha confuso molti studenti: la
        \emph{challenge} è contro una sfida, la \emph{competition} è contro
        altre persone, altri giocatori.
  \item[Accomplishment (realizzazione, progressione)] salire di livello,
        sbloccare la nuova armatura più potente, raggiungere il ramo
        dell'albero delle abilità che si desiderava: la sensazione di crescere
        dà senso a ciò che si fa.
  \item[Socialization (socializzazione)] parlare con la gente: l'essere umano
        è un animale sociale, e interagire con le persone può essere una fonte
        di gioia (e non solo).
  \item[Narrative (narrativa)] le storie: una storia appassionante, una
        storia triste, una storia comica, una battuta.
  \item[Informative (apprendimento)] imparare. L'anime \emph{Dr.\ Stone},
        tra i preferiti del direttore, gli piace proprio perché è informativo:
        insegna qualcosa di vero.
  \item[Aesthetic (estetica)] le cose \emph{visivamente piacevoli}. Esistono
        regole, come quelle sui codici colore, che indicano combinazioni di
        colori e forme generalmente considerate piacevoli, cioè non fastidiose
        alla vista. Il ``raggio epilessia'' è l'opposto.
  \item[Auditory/visual stimulation (stimolazione uditiva e visiva)] lo stesso
        discorso per i suoni. La sirena (creata dall'uomo) e il pianto di un
        bambino (creato dalla natura) sono fatti per essere fastidiosi, per
        spingere ad agire: la sirena per far spostare, il pianto perché il
        bambino, che non sa fare nulla, venga accudito. All'opposto, esistono
        suoni con caratteristiche piacevoli all'orecchio.
  \item[Humor (umorismo)] è separato dalla narrativa, di cui potrebbe far
        parte, perché far ridere è difficilissimo: l'umorismo è una delle cose
        più personali in assoluto, e quello di una persona può essere
        completamente diverso da quello di un'altra. Riuscire a fare un
        prodotto umoristico è difficilissimo.
  \item[Physical manipulation (manipolazione fisica)] il direttore ha giocato
        tutta la lezione con un fazzolettino, perché è piacevole: la texture,
        il poterlo piegare. I fidget spinner si basano su questo: usare le
        mani è piacevole.
  \item[Physicality (fisicità)] lo stesso discorso ma con il corpo: correre con
        le scarpe giuste, fare l'attività fisica che piace.
  \item[Creation (creazione)] la creatività. Gli studenti sono lì perché sono
        creativi, chi più chi meno, chi in un modo chi in un altro: hanno il
        desiderio di aggiungere al mondo qualcosa che prima non c'era (si
        pensi a chi ha ricostruito Westeros in Minecraft).
  \item[Vicarious experience (esperienza vicaria)] fare cose che non sono
        comuni o che non esistono: il direttore non è un pilota di Formula 1,
        non vuole fare il soldato, non è uno stregone che combatte i draghi, ma
        può trovare piacevole vivere queste esperienze in forma simulata. Ci
        sono vari livelli: da ciò che non possiamo fare ma un essere umano
        potrebbe fare, a ciò che è del tutto al di fuori delle nostre
        possibilità, come essere un alieno in una galassia lontana lontana.
  \item[Logic/Puzzle (logica, problem solving)] usare la propria mente per
        risolvere un problema.
\end{description}
Da subito bisogna usare questi termini; pian piano il vocabolario si
espanderà.

% ---------------------------------------------------------------------
\section{La forma mentis del designer secondo Chris Bateman}
\label{sec:bateman}

L'ultima parte della lezione riprende alcuni consigli di Chris Bateman
(autore, tra l'altro, di \emph{21st Century Game Design}, si veda la
sezione~\ref{sec:risorse}) per acquisire la \emph{forma mentis} del designer.

\begin{description}
  \item[Il game design è una metodologia iterativa ed evolutiva] davanti a un
        problema bisogna procedere in modo \textbf{scientifico}. Non si può
        realizzare semplicemente la prima idea che viene in mente: bisogna
        generare un gran numero di idee, sia con la creatività sia con le
        metodologie (come i framework di analisi), e analizzare il problema da
        tutti i punti di vista, per trovare non \emph{una} soluzione, ma
        \textbf{la soluzione migliore}. Una soluzione la trovano tutti: per
        risolvere il problema delle tasse in Italia si potrebbero, per
        paradosso, eliminare tutti i contribuenti; è una soluzione, ma non è
        certo la migliore. La soluzione migliore va analizzata e trovata.
  \item[Non esiste un unico metodo di design] è quanto detto a proposito
        delle lenti. Più metodi si studiano e più tecnologie si imparano, più la
        mente si espande e si trovano modi di interagire con problemi, persone e
        design, e si diventa designer migliori. \emph{The right tool for the
        right job}; ma uno strumento non basta, ne servono centomila.
  \item[Il design è \emph{egoless}] durante la lezione i docenti hanno
        smontato le definizioni degli studenti senza avercela minimamente con
        loro, e gli studenti avrebbero potuto smontare la loro: si stava
        discutendo di teoria del design. \textbf{Le idee non rappresentano chi
        le ha}: sono un elemento separato, da analizzare, distruggere,
        ricostruire, migliorare. Che un designer generi cento idee pessime non
        significa che sia incompetente: le idee vanno generate e analizzate.
        \textbf{Le idee possono venire da chiunque}: in uno dei giochi del
        direttore, una delle meccaniche migliori è stata suggerita da un
        programmatore. Il design consiste nell'analizzare l'idea, comprenderla,
        migliorarla e capire se funziona. Altre volte è stato Mattia, o Jack,
        lavorando con lui, a dirgli ``sei sicuro che sia così?'', e avevano
        ragione; questo non ha minimamente intaccato il suo valore di designer.
        Si crea, si analizza e si distrugge, sempre con lo scopo di trovare la
        soluzione migliore, qualunque essa sia.
  \item[Il design è guidato dagli utenti] lo scopo finale, come molte
        definizioni avevano colto, è creare un gioco \textbf{che venga
        giocato}, che l'utente finale apprezzi. Non un gioco che il designer
        considera il migliore del mondo e che tutti gli altri odiano. Fare arte
        per se stessi è un approccio del tutto accettabile, e non c'è nulla di
        male nell'arte fine a se stessa: anche il direttore fa giochi ``per
        Diego'', esperimenti e tentativi, e ci saranno momenti per farli (ad
        esempio nelle \emph{game jam}, che verranno spiegate più avanti). Ma il
        corso insegna a lavorare come designer, e il game design come
        professione è orientato all'utente finale: creare un prodotto che una,
        si spera decine o centinaia di milioni di persone possano giocare e
        apprezzare (il che, detto brutalmente, se succede rende anche molto
        ricchi). Il direttore, che fa consulenza a vari studi, ha visto spesso
        team fare un gioco, vederlo fallire e perdere tutti il lavoro: un
        peccato, perché con lo stesso studio si sarebbero potuti creare altri
        capolavori, o semplicemente dare lavoro alle persone.
\end{description}

% ---------------------------------------------------------------------
\section{Il task della settimana}
\label{sec:task-gd2}

\begin{enumerate}
  \item Scegliere un \textbf{gioco arcade}.
  \item Scegliere \textbf{uno} dei framework visti a lezione: Schell, Romero e
        Schreiber, oppure Bura (quest'ultimo è più complesso, ma accettato).
  \item Analizzare il gioco usando il framework scelto.
  \item Far emergere le \textbf{emozioni} verso cui il gioco è spiccatamente
        orientato, usando il vocabolario visto in lezione: individuarne
        principalmente \textbf{tre}.
  \item Poiché in aula è stato chiesto che cosa si intenda per ``gioco
        arcade'', aggiungere al task anche \textbf{una propria definizione di
        gioco arcade}.
\end{enumerate}
Vincoli e indicazioni:
\begin{itemize}
  \item valgono tutti i \textbf{termini banditi} (sezione~\ref{sec:vocabolario});
  \item le modalità di consegna (ad esempio documento su Drive) e le date
        verranno comunicate da ciascun docente per la propria sede; come regola
        generale, un task si considera da consegnare entro una settimana
        dall'assegnazione;
  \item lo studente che ha pronunciato la parola proibita deve svolgere
        l'analisi con \textbf{entrambi} i framework principali, in
        \textbf{due documenti distinti};
  \item la lezione successiva si aprirà con una sessione di domande e risposte
        su questa lezione.
\end{itemize}

% ---------------------------------------------------------------------
\section*{Punti chiave del capitolo}
\addcontentsline{toc}{section}{Punti chiave del capitolo}

\begin{itemize}
  \item Il primo passo del percorso è passare dalla visione del
        \textbf{giocatore} a quella del \textbf{creatore}.
  \item Ogni definizione di gioco include qualcosa ed esclude qualcos'altro: le
        parole hanno un significato preciso e vanno scelte con cura.
  \item Il gioco \textbf{non è definibile}, perché è anche uno \textbf{stato
        mentale}; \emph{game} e \emph{play} sono la stessa cosa, e lavorare con
        lo stato mentale del giocatore è parte del mestiere.
  \item La propria definizione di gioco è una \textbf{bussola}: dice che tipo
        di designer si è. Va riscritta periodicamente, senza rileggere la
        precedente.
  \item Non si può definire la zuppa, ma se ne possono elencare gli
        \textbf{ingredienti}: definire chiude, elencare apre.
  \item Un framework valido deve potersi applicare a \textbf{qualsiasi gioco}.
  \item \textbf{Schell}: meccaniche, narrativa, estetiche, tecnologia.
        \textbf{Romero e Schreiber}: game qualities, giocatore, tema,
        obiettivo, regole, risorse, interazione. \textbf{Bura}: action, system,
        self, social incrociati con freedom, mastery, data.
  \item Non esiste un metodo unico: i framework sono \textbf{lenti}, e servono
        a guardare un problema da più punti di vista. È anche la miglior difesa
        del designer contro l'automazione.
  \item La parola ``divertente'' (e con lei ``bello'', ``noioso'', ecc.) è
        \textbf{bandita}: bisogna nominare le \textbf{emozioni specifiche} che
        il gioco vuole suscitare.
  \item Il design è iterativo, plurale, \textbf{egoless} e orientato
        all'\textbf{utente}.
\end{itemize}
